% SOURCE docs/release/2026-10-29/source/docs/writing-repair/2026-10-29/src/ba7c02f8904143a4.tex CHAR 0
\chapter{Keeping the economic question intact across methods}
\label{in:semantics}

An analysis remains the same analysis only if its population, intervention, information and uncertainty survive each computational step. This chapter first develops that requirement for coordinates, graph classes and sets. It then examines the evidence from local custody, recovery and interface trials, before comparing restricted Python identification, Java discovery and R/native forest routes. Each route has a different scientific domain. Their common representation makes those domains visible; it does not make their guarantees interchangeable.

\section{The unit of exchange is an analysis with meaning}

A funding service, a servicing notebook and an analyst's agent may all request
an ``effect'', while referring to different populations, actions, horizons and
units. Exchanging a floating-point number is therefore insufficient. The
proposed common object contains a scientific context, an operation, input
identities, an execution disposition and a qualified result. The local
reference implementation enforces selected parts of this contract; it does not
validate the truth of every declaration. \citep{v12_in_contract,v12_in_runtime}

For a scalar intervention query, write the context schematically as
\[
 Q=(P,\mu,a,a_0,Y,H,u,s,\mathcal T,\mathcal A,G),
\]
where $P$ is the population, $\mu$ its weighting measure, $a$ and $a_0$ the
intervention and comparator, $Y$ the outcome, $H$ the horizon, $u$ the units,
$s$ the sign convention, $\mathcal T$ the information clocks, $\mathcal A$ the
assumptions and sampling declarations, and $G$ the graph or graph class.
This mathematical notation expresses the required meaning; particular schema
fields may encode only a supported subset. The original declaration must be
retained when a narrower operation representation is used.

Changing $a-a_0$ from a ten-basis-point offer to a one-percentage-point offer
changes the question, even if the same coefficient field can store both.
Replacing customer weights by balance weights changes $\mu$. Treating a
post-decision observation as available history changes $\mathcal T$. A digest
can detect that the declared context changed. It cannot establish that the
context describes the actual banking process.

The implementation binds supplied moments and sensitivity artifacts to their
scientific context. Reusing their bytes with a changed sign, clock, unit or
target is rejected unless a new, appropriate artifact is supplied. This
prevents a consumer from reinterpreting a cached numerical object through an
unannounced schema change. It does not replace a scientific argument for the
moments or covariance. \citep{v12_in_runtime}

\section{Execution, science and qualification are different states}

A successful invocation can return unsupported identification, an empty
sensitivity set, or a whole-line confidence set. These are not necessarily
transport failures. Conversely, a timeout supplies no scientific answer even
if the requested mathematical object exists. The common response separates
invocation state, scientific disposition and qualification evidence.

This separation gives a batch operator a defensible continuation rule. A
completed negative answer can be archived and used in a report. A malformed
request needs correction. An unsupported operation needs another qualified
backend or a narrower, explicitly authorized question. A contradiction in the
declared premises needs scientific review. Automatically retrying all four
events until a finite number appears would select answers according to their
convenience.

The library's supplied-DAG back-door operation illustrates the distinction.
It can attempt an adjustment expression in its bounded domain. Failure of
that attempt does not prove general nonidentification; a different graphical
functional or additional information could be relevant. The response must
describe the actual method and its limitation. Likewise, graph syntax review
checks declared structures but does not discover their causal truth.
\citep{v12_in_runtime}

\section{Sets must retain their topology}

The local representation distinguishes unknown from empty, and permits a union
of intervals with infinite endpoints. Unknown has no asserted interval list;
an empty set has an empty list. A whole-line set has one interval unbounded in
both directions. A disconnected set keeps its components. \citep{v12_in_contract}
None should be represented by an arbitrary plotting range.

Consider the stipulated effect set
\[
 \Theta=(-\infty,-2]\ \cup\ [3,\infty)
\]
and a known monetary map $V=100\theta-300$. The image is
\[
 V(\Theta)=(-\infty,-500]\ \cup\ [0,\infty).
\]
Taking the convex hull would return the entire real line and discard a real
exclusion. Clipping the tails would falsely assert boundedness. Changing to a
loss convention $L=-V$ gives $(-\infty,0]\cup[500,\infty)$; endpoint order must
reverse when multiplying by a negative number.

If $\Theta$ has a stated coverage guarantee for $\theta$, applying a known
deterministic map preserves containment on the same event. It does not create
a guarantee for an estimated conversion factor or a randomly chosen target.
When margins, discount rates or exposure weights are estimated, joint
uncertainty or an explicitly conditional analysis is needed. Serialization
cannot make uncertain financial inputs fixed.

The same care applies to graph sets. Member identities, endpoint semantics and
evidence must travel with the alternatives. A CPDAG, PAG and selected DAG are
not interchangeable edge lists. The present schema retains declared classes
and hashed membership; its validator is not a complete recognizer or confidence
construction for arbitrary temporal equivalence classes.

\section{A contract for future numerical adapters}

Each adapter should state its accepted targets, graph classes, sampling laws,
missingness rules, numerical inputs and output uncertainty. A source package
may use an ATT where the caller requested an ATE, return pointwise intervals
where simultaneous inference is required, or approximate a set with a hull.
Those are semantic mismatches even when both sides exchange valid JSON.

The adapter must either preserve the requested meaning, return a disclosed
qualified approximation, or refuse the mapping. A proposed R, Python, Java or
C++ bridge therefore needs matched fixtures at ordinary and adverse boundaries,
including disconnected sets, weak identification and unsupported models.
Cross-language execution, licenses, data copying and complete workloads remain
qualification obligations. The existing local dispatcher is a common boundary
for that work, not evidence that those bridges already exist.


\hypertarget{when-a-conversion-preserves-the-question}{%
\section{When a conversion preserves the
question}\label{when-a-conversion-preserves-the-question}}

The interchange obligation can be stated as a commutation requirement.
Let \(F_Q(D)\) be a scientific calculation on observations \(D\) for
context \(Q\). Let \(E\) encode the observations and question for an
engine, \(M\) denote that engine's computation, and \(R\) decode its
result. If \(T\) is a declared change of coordinates, an exact adapter
should satisfy

\[
 R\{M(E(D,Q))\}=T\{F_Q(D)\}.
\]

The right side specifies what the answer means before serialization. The
left side describes the proposed software path. This equality is an
obligation on the admitted domain, not a theorem established by two
programs returning parseable JSON. It includes a negative scientific
answer or a set with several components whenever that is the defined
output of \(F_Q\). A method that changes the target while improving
numerical stability has changed \(Q\); it has not merely implemented
\(E\).

There are three useful levels of equivalence. Exact structural
equivalence preserves a graph, variable map or expression tree up to
declared renaming. Numerical equivalence permits a quantified error in
the requested coordinates. Scientific equivalence additionally requires
the same population, intervention, sampling law and uncertainty
interpretation. A fixture showing agreement of two printed means
addresses only a small part of the second level. The design should state
which level it establishes and where further evidence is required.

An adapter also needs an explicit domain. Write \(\mathcal D_M\) for the
queries and inputs admitted by engine \(M\). If
\((D,Q)\notin\mathcal D_M\), the proper response is a typed unsupported
result that preserves the original request. Automatically dropping
weights, selecting one graph, clipping an interval or replacing a
temporal sample with independent rows does not extend the domain
honestly. A proposed approximation can still be useful, but its
transformed question and additional assumptions must be disclosed as a
different operation.

This view explains why a generic language bridge is insufficient. A
Python-to-R array conversion can preserve all floating-point values
while losing the fact that columns represent treatment arms in a
specific order. A Java graph object can preserve adjacency while
reversing an endpoint mark. A C++ routine can accurately optimize an
objective for a different weighting measure. The mathematical interface
must bind these meanings alongside the numeric payload, and its review
must include examples for which apparently harmless changes lead to
materially different bank decisions.

\hypertarget{population-measures-and-rate-coordinates}{%
\section{Population measures and rate
coordinates}\label{population-measures-and-rate-coordinates}}

Consider a bank's effect target

\[
 \tau_\mu=\int\{m_1(x)-m_0(x)\}\,d\mu(x),
\]

where \(m_a(x)\) is a conditional mean potential outcome under action
\(a\), and \(\mu\) is the declared distribution of customers or
exposures. The same response functions can imply different targets under
different measures. In an assumed two-segment example, the customer
shares are \(3/4\) and \(1/4\), the segment effects are USD 20 and USD
-40 per account, and relative balances are one and nine. The
customer-average effect is

\[
 \tfrac34(20)+\tfrac14(-40)=5.
\]

Balance weighting gives masses \(1/4\) and \(3/4\), and hence an effect
of USD -25. The sign reverses even though no within-segment estimate
changed. The example is stipulated arithmetic, not evidence about the
distribution of any bank's customers. Its purpose is to make the
weighting measure an indispensable part of the requested estimand.
\citep{exec13_interchange_examples}

A caller may intentionally prefer the balance-weighted target for a
treasury decision. That choice should create a new target identity and
explain which decision it serves. It cannot be inferred from a column
named \texttt{weight}, because a supplied weight might instead describe
unequal sampling probabilities, treatment assignment, transport to
another population or a numerical objective. These roles have different
identification and variance implications. Multiplying several weights
without preserving their definitions can conceal that distinction.

Treatment coordinates require similar care. Let \(r\) be a decimal
interest rate, so a basis-point change is \(10^{-4}\). If a linear
response derivative is \(\beta=\partial E(Y)/\partial r\), then the
derivative per basis point is \(10^{-4}\beta\). A hypothetical estimate
of USD 2 million per unit decimal rate becomes USD 200 per basis point;
an interval of USD 1--3 million becomes USD 100--300. Reporting the
original number under the new label would create a factor-of-10,000
error. The exact calculation is reproducible. It does not assert that a
constant derivative is an appropriate response model over an arbitrary
offer range. \citep{exec13_interchange_examples}

For a finite nonlinear rate change, the target is instead
\(E\{Y(r_1)-Y(r_0)\}\); rescaling a local derivative does not identify
that contrast. Nor does a per-month response become an annual causal
response by changing its unit label. Timing, compounding, feedback and
the intervention path may matter. Units should therefore describe the
mathematical coordinate actually estimated, while the horizon and
intervention path remain separate fields. A conversion that requires new
structural assumptions belongs in a scientific transformation, not a
formatting function.

\hypertarget{joint-uncertainty-across-coordinates}{%
\section{Joint uncertainty across
coordinates}\label{joint-uncertainty-across-coordinates}}

Suppose an admitted calculation returns \(\widehat\theta\) and
covariance \(\widehat\Sigma\) in the declared coordinates. Under a fixed
linear map \(\phi=L\theta+c\), the corresponding estimate and covariance
are

\[
 \widehat\phi=L\widehat\theta+c,
 \qquad \widehat\Sigma_\phi=L\widehat\Sigma L^\top.
\]

This covariance identity follows by expanding the centered vector
\(L(\widehat\theta-E\widehat\theta)\). It preserves whatever
justification the original covariance actually has; it supplies no new
consistency, normal approximation or finite-sample coverage argument. If
\(L\) is estimated, choosing to hold it fixed is a conditional analysis
and must be labeled. Joint estimation error cannot be eliminated by
storing the fitted \(L\) as a configuration value.

For an exact arithmetic illustration, let the first coordinate be a
gross revenue effect in USD and the second an operating-cost effect in
thousands of USD. Suppose the supplied covariance is

\[
 \Sigma=\begin{pmatrix}40000&20\\20&1/4\end{pmatrix},
 \qquad L=\begin{pmatrix}1/1000&0\\0&-1\end{pmatrix}.
\]

The new coordinates are gross benefit and negative operating cost, both
in thousands of USD. Their covariance is

\[
 L\Sigma L^\top=
 \begin{pmatrix}1/25&-1/50\\-1/50&1/4\end{pmatrix}.
\]

The negative cross-covariance matters when adding the two coordinates to
obtain net value: its variance is \(1/4\), whereas dropping that
covariance gives \(29/100\). A column-order error or a failure to
reverse the covariance sign changes the uncertainty even if both
marginal standard errors look reasonable. The matrices here are assumed,
positive definite reference objects whose conversion is checked by exact
arithmetic. They are not fitted bank uncertainty estimates.
\citep{exec13_interchange_examples}

Marginal intervals cannot substitute for a joint region. If \(C_1,C_2\)
each cover their own fixed coordinate with probability at least \(.95\),
the union bound gives only

\[
 P\{(\theta_1,\theta_2)\in C_1\times C_2\}\geq .90.
\]

Even with independent coverage events and exact marginal coverage, the
joint probability is \(.95^2=.9025\). A simultaneous 95\% statement
needs a construction that supports it; a transport cannot supply that
construction by changing the uncertainty label. The same issue arises if
an agent inspects several analyses and reports only the most favorable
one. Coverage for each fixed query is not automatically coverage for the
selected query. Preserving the full query-selection rule and any
simultaneous or selective argument is part of the workflow's scientific
identity.

These distinctions already appear in the local semantic validator:
confidence, posterior credibility and sensitivity have separate
meanings, a joint-confidence declaration requires construction evidence,
and the declared sampling law must match the context. Those checks
constrain what can be serialized. They do not independently prove the
submitted construction or establish Bayesian calibration.
\citep{exec13_interchange_runtime}

\hypertarget{a-confidence-set-is-not-a-pair-of-convenient-endpoints}{%
\section{A confidence set is not a pair of convenient
endpoints}\label{a-confidence-set-is-not-a-pair-of-convenient-endpoints}}

Weak identification supplies a practical reason to retain topology. For
a ratio target \(\theta=\mu_Y/\mu_X\), one can invert a test of
\(\mu_Y-\theta\mu_X=0\). With supplied approximately Gaussian mean
errors, write their variances as \(s_{YY},s_{XX}\) and covariance as
\(s_{XY}\). A residual-test set takes the form

\[
 (\bar Y-\theta\bar X)^2
 \leq c^2(s_{YY}-2\theta s_{XY}+\theta^2s_{XX}).
\]

Its interpretation depends on the stated Gaussian approximation,
covariance and critical value; algebraic inversion alone is not a
finite-sample theorem for an arbitrary dependent sample. Collecting
terms gives

\[
 a\theta^2+b\theta+d\leq0,
\]

where

\[
 a=\bar X^2-c^2s_{XX},\quad
 b=-2\bar X\bar Y+2c^2s_{XY},\quad
 d=\bar Y^2-c^2s_{YY}.
\]

When the denominator is weak, \(a\) may be negative. With two real roots
the solution then consists of exterior rays; with other coefficients it
may be the whole line or another topology. The point ratio can fail to
exist while the test-inversion set still has a precise meaning. For the
stipulated values \(\bar Y=2,\bar X=0,c=2,s_{YY}=s_{XY}=0,s_{XX}=1/4\),
the inequality is \(4\leq\theta^2\), giving

\[
 (-\infty,-2]\cup[2,\infty).
\]

This degenerate but positive-semidefinite reference covariance makes the
algebra transparent. It does not validate that covariance for a
data-generating process. It shows why a response object must allow a
valid set without a finite point ratio. The existing local
supplied-moment operation retains these distinctions under its declared
approximation.
\citep{exec13_interchange_examples,exec13_interchange_runtime}

Unknown, empty, whole-line and disconnected results answer different
questions. Unknown says that no set has been computed. Empty may signal
incompatible restrictions in a particular construction, and does not
mean an effect of zero. Whole-line uncertainty means every real value
remains allowed by that construction. An exterior-ray set excludes a
middle region despite being unbounded. A graph-class result may instead
be a set of distinct functionals whose numerical evaluation is not yet
available. None can be reduced to a nullable scalar without losing
information that could change a decision.

Nonlinear monetary transformations make the loss visible. Take the
assumed effect set \([-3,-2]\cup[2,3]\) and value map
\(V(\theta)=100(\theta^2-1)\). Its image is \([300,800]\). Applying the
same map to the convex hull \([-3,3]\) gives \([-100,800]\). Both are
outer possibilities for different retained information, but the latter
invents a negative-value possibility that the original set excludes. A
hull may be a disclosed conservative approximation; it must not be
described as exact interchange. The reverse mistake, returning only
transformed outer endpoints without considering interior extrema, can
omit valid values and lose coverage. \citep{exec13_interchange_examples}

\hypertarget{combining-statistical-measurement-and-numerical-uncertainty}{%
\section{Combining statistical, measurement and numerical
uncertainty}\label{combining-statistical-measurement-and-numerical-uncertainty}}

Suppose an effect set \(C\) contains the fixed target \(\theta\) with
probability at least \(1-\alpha_C\), and a set \(H\) contains a needed
valuation parameter \(\eta\) with probability at least \(1-\alpha_H\).
For a declared value function \(v\), the projection

\[
 W=\{v(t,h):t\in C,h\in H\}
\]

contains \(v(\theta,\eta)\) whenever both input events hold. Thus the
union bound yields coverage at least \(1-\alpha_C-\alpha_H\), without
independence. A valid joint region can be projected instead and may be
narrower because it retains dependence between \(\theta\) and \(\eta\).
Independently varying marginal ranges is generally an outer
approximation to that joint region. Fixing \(\eta\) at an estimate has
no such unconditional guarantee without an additional argument.

Numerical containment is a third obligation. Let \(\widetilde W\) be an
approximate computed set. If a certificate establishes

\[
 W\subseteq\widetilde W\oplus[-\varepsilon,\varepsilon],
\]

where \(\oplus\) denotes addition of sets, then the enlarged computed
set preserves the preceding coverage event. If the numerical certificate
itself can fail with probability \(\delta\), the resulting guarantee is
at least \(1-\alpha_C-\alpha_H-\delta\). Budgets of \(.02,.02,.01\) give
a 95\% lower guarantee. This arithmetic does not prove that a particular
optimizer supplies the assumed certificate; its error bound must come
from the numerical method. \citep{exec13_interchange_examples}

A grid of feasible points usually supplies an inner sample, not the
required outer enclosure. Missing a disconnected component or an
unbounded tail is therefore a scientific failure, even if all sampled
points satisfy the constraints. An optimizer's small objective change
also need not certify the set boundary. For a nonlinear map, a Lipschitz
constant can propagate an input approximation error only over a region
where the constant is valid; a bound obtained on a clipped plotting
window does not control unbounded tails outside it.

The currency of the numerical error should be the currency of the
decision. If a statistical lower value is USD 30 but an admitted
numerical approximation has error up to USD 40, the certified lower
value is USD -10. Reporting only the statistical endpoint can reverse a
claimed robust recommendation. Conversely, a deliberately coarse
computation may be adequate when its certified error cannot change the
decision. This makes numerical precision an explicit investment tradeoff
rather than an unexamined default tolerance.

\hypertarget{sampling-declarations-and-partial-scientific-answers}{%
\section{Sampling declarations and partial scientific
answers}\label{sampling-declarations-and-partial-scientific-answers}}

A shared data format does not create shared sampling assumptions.
Consider the generated one-date account model \(Y_i=\theta+C+E_i\),
where independent centered Gaussian components have
\(\operatorname{Var}(C)=.09\), \(\operatorname{Var}(E_i)=1\), and \(C\)
is shared by every account. Then

\[
 \operatorname{Var}(\bar Y)=.09+1/N.
\]

At \(N=10{,}000\), the correct variance is \(.0901\). Treating the rows
as independent with their individual variance \(1.09\) gives
\(.000109\), understating variance by \(90100/109\), or about 827 times.
The number of rows was preserved exactly; the scientific information was
not. This exact generated calculation illustrates why a panel or
calendar declaration must survive conversion and why an iid-only method
must reject it. It makes no claim about an observed bank covariance.
\citep{exec13_interchange_examples}

The same problem occurs with persistent linkage errors, cluster
treatment assignment and policies that condition on shared history. A
consumer that sees only row count and a floating-point estimate cannot
reconstruct the lost sampling unit. The local bounded back-door
reference explicitly restricts its sampling regime, while the proposed
symbolic identification adapter preserves the declaration without
supplying estimation uncertainty. Those are different services with
different obligations.
\citep{exec13_interchange_runtime,exec13_y0_contract}

Finally, the response should identify what has actually been
established. A source-inspected method may support a proposal decision.
An executed adapter may establish that particular conversions work. A
statistical region may be valid under declared scientific assumptions. A
banking action may still be undecided because its loss function or
information is insufficient. Preserving a partial answer makes the next
step precise: obtain a needed assumption, improve information, choose
another method or refine numerical containment. Automatically converting
every partial answer into a confident recommendation would obscure the
very gap the product is meant to explain.

\hypertarget{preserving-causal-meaning-in-a-python-identification-adapter}{%
\section{Preserving causal meaning in a Python identification
adapter}\label{preserving-causal-meaning-in-a-python-identification-adapter}}

\label{exec13:python-identification}

A common causal interface is useful only if a question keeps its meaning
when it passes into another library. An outreach decision can become a
different question if a transport drops a condition, converts an
uncertain graph into one directed graph, or treats a symbolic
probability formula as an estimated monetary return. The proposed Python
path therefore begins with a narrow service: identify an intervention
distribution in one supplied causal graph, preserve the assumptions that
make the formula relevant, and report an explicit limit when that
service cannot answer the question.

y0 is a plausible component for this service because its inspected
Python API separates a query, a mixed graph and a symbolic
identification expression. The source considered here is commit
\texttt{cc6644d3b470a1ab2a1eb9da4f0fdb0c843f5d3a}, dated 13 April 2026;
its project metadata declares development version \texttt{0.2.12-dev}.
The acquisition is later than the literature cutoff, but the reviewed
commit predates 7 October 2026. These identifiers describe inspected
source, not a tested deployment. The proposed decision is to integrate a
restricted identification path and verify its conversions before
admitting it to a prototype. Porting an entire identification engine
would create a second implementation to maintain before the first
interface has established its scientific meaning.
\citep{exec13_y0_api,exec13_y0_metadata}

\hypertarget{a-distributional-question-before-an-estimate}{%
\subsection{A distributional question before an
estimate}\label{a-distributional-question-before-an-estimate}}

Let \(A\) denote a bank action, \(Y\) an outcome and \(Z\) a collection
of explicitly named conditions. The requested object is

\[
 p(y\mid\operatorname{do}(a),z).
\]

It is distinct from an observed conditional probability, an average
difference between two actions, and a confidence interval for an
estimated difference. For a finite outcome support, a later evaluator
might form

\[
 \mu(a,z)=\sum_y y\,p(y\mid\operatorname{do}(a),z),
 \qquad \tau(z)=\mu(a_1,z)-\mu(a_0,z).
\]

That calculation requires the original outcome coordinates,
treatment/control order and units. If \(Y\) is a repayment indicator,
\(\tau\) is a probability difference. If \(Y\) is net cash received, it
has the declared currency and horizon. Multiplying the first by a loss
avoided per repayment is an additional valuation model; it is not an
operation that the identification library can infer from a variable
label.

The first adapter is deliberately limited to one scalar outcome and
treatment, optional conditions, one causal DAG or acyclic directed mixed
graph (ADMG), and declared finite discrete domains. This makes its
numerical references and interpretation concrete. It does not establish
an absence of wider functions elsewhere in y0. A graph representing
selection, transport between populations, cross-world counterfactuals,
an adaptive longitudinal regime or unresolved graph alternatives needs a
separately reviewed path. An account panel can still describe its
sampling law in the request, but a symbolic formula supplies no
independence assumption, standard error or calibration for that panel.

The scientific context should travel with the formula: target
population, action values, outcome horizon, observation and availability
clocks, variable domains, units, sign convention, sampling unit and
assumptions. A hash of that complete context identifies the request. A
changed population, a reversed action contrast or a different
availability clock produces a different identity even when the graph
picture is unchanged. The shared library, CLI and MCP operation
contracts already distinguish transport completion from scientific
disposition; the proposed adapter specializes that distinction for
identification. \citep{exec13_y0_contract}

A condition deserves particular care. Conditioning on a pretreatment
risk category and conditioning on a response to the action are different
targets. The interface should require the caller to state which is
intended, preserve the chronology, and avoid silently relabeling a
post-action variable as an adjustment covariate. Syntax can reject a
missing clock declaration. It cannot determine whether the bank's
claimed chronology or causal model is true.

\hypertarget{graph-and-coordinate-conversion}{%
\subsection{Graph and coordinate
conversion}\label{graph-and-coordinate-conversion}}

The graph object in the inspected source contains a directed component
and a component named \texttt{undirected}. In the causal mixed-graph
interpretation, the latter carries bidirected latent links: the source's
causal-graph importer maps the input type \texttt{bidirected} into that
field, and its latent expansion creates a shared hidden parent for each
such edge. It is therefore incorrect to put an unresolved CPDAG edge
into this component merely because both are stored by an undirected
graph data structure. \citep{exec13_y0_graph}

For the admitted path, a tail--arrow pair becomes a directed edge, and
an arrow--arrow pair becomes an ADMG bidirected edge. Circles,
tail--tail edges and graph collections receive an unsupported
disposition that retains the original representation. A directed edge
and a bidirected edge may coexist on the same pair of variables; erasing
one would change the model. The directed component must be acyclic, with
all endpoints in the declared node set. These checks belong in the
adapter: the inspected graph initializer sets metadata flags and does
not itself establish acyclicity.

Every variable, including an isolated one, must be passed explicitly.
Constructing a graph solely from edges can lose an outcome or an
independent condition that has no incident edge. The adapter should
compare the node and typed-edge sets before and after conversion and use
a deterministic valid topological ordering. The purpose is
reproducibility and semantic preservation, not the selection of a
statistically preferred causal ordering.

Names are another coordinate system. External identifiers may contain
spaces, non-ASCII text or labels that collide after informal
sanitization. The proposed conversion assigns opaque internal symbols
through a bijection and keeps a sidecar containing the external
identifier, displayed label, domain order, time coordinate and units. It
never guesses temporal order from a name such as \texttt{lag1}, and it
never converts category labels into numeric doses without a declared
mapping. A graph-set request remains a graph-set request even when
unsupported. An uncomputed effect set remains unknown; it is not
represented by a convenient point estimate or an empty region.

\hypertarget{why-conditional-identification-needs-a-restricted-route}{%
\subsection{Why conditional identification needs a restricted
route}\label{why-conditional-identification-needs-a-restricted-route}}

The high-level API calls ID when its conditions field is empty and IDC
otherwise. By default it converts the specific \texttt{Unidentifiable}
exception to \texttt{None}; with \texttt{strict=True} it preserves the
exception. Other exception types are not caught by that branch. The
adapter should preserve this distinction, because a malformed graph, a
runtime failure and nonidentification are different outcomes. The
inspected ID recursion supplies node/district arguments when raising
\texttt{Unidentifiable}, but those arguments should not be presented as
an independently verified graphical certificate.
\citep{exec13_y0_api,exec13_y0_id}

Two further source findings restrict this conditional route. First, the
\texttt{Query.expression} property builds a conditioned distribution when
conditions exist. Its intervention path is an \texttt{elif} branch. When
both are present, that convenience property does not add the
intervention. This is a serialization hazard, not evidence that the
separate IDC algorithm discards its treatment field:
\texttt{Identification} retains outcomes, treatments and conditions
separately. The adapter should serialize its original query triple
independently of that convenience property. \citep{exec13_y0_utils}

Second, IDC tests whether an observation may be exchanged for an action
using a separation routine. The inspected routine takes a
directed-ancestor subgraph, connects directed co-parents, flattens the
remaining edges and removes conditioning vertices. For a DAG this is a
familiar construction. The same steps do not generally implement the
collider behavior of bidirected paths. The source-level discrepancy can
be demonstrated without relying on a package's own numerical output.
\citep{exec13_y0_idc,exec13_y0_rule2,exec13_y0_separation,exec13_y0_graph}

Consider the ADMG

\[
 Y\leftrightarrow B\leftrightarrow Z,
\]

with an unrelated action \(X\). The middle vertex \(B\) is a collider.
Conditioning on it can connect \(Y\) and \(Z\); treating it as an
ordinary undirected intermediate vertex and deleting it instead
disconnects them. In the rule-2 check for exchanging \(Z\), removal of
incoming edges to isolated \(X\) and outgoing edges from \(Z\) does not
change this path. The inspected directed co-parent rule adds no
\(Y\)--\(Z\) link because there are no directed edges. Its subsequent
deletion of \(B\) therefore predicts separation where the causal mixed
graph does not.

An exact generated law makes the consequence explicit. Let \(U,V\) be
independent fair bits, let \(Y=U\), \(Z=V\), and let \(B\) report
whether \(U=V\), correctly with probability \(.9\). The observed law of
\((Y,B,Z)\) has positive probability in all eight cells. Direct
enumeration gives

\[
 P(Y=1\mid B=1,Z=1)=\frac9{10},
 \qquad
 P(Y=1\mid B=1,\operatorname{do}(Z=1))=\frac12.
\]

An intervention on \(Z\) overwrites that measurement but does not change
\(U,V\) or the mechanism generating \(B\). Exchanging the observation
for the action would therefore change the requested probability. This is
an exact reference for the scientific meaning of the graph and a trace
of the inspected source path; it is not a report of an executed y0
failure. \citep{exec13_y0_oracles}

The proposed adapter can avoid that native conditional path, subject to the internal node-preservation qualification in Section~\ref{exec13:isolated-condition}. Its mathematical route first applies ID to the joint query \(p(y,z\mid\operatorname{do}(a))\). If it
obtains an observational functional \(N(y,z,a)\), it forms

\[
 \frac{N(y,z,a)}{\sum_{y'}N(y',z,a)},
\]

on coordinates for which the denominator is positive. This is a direct
probability identity. It neither uses the questionable separation step
nor relies on the query-expression convenience property. The route is
conservative: failure to identify the joint distribution does not prove
failure to identify the conditional target. Such a failure must return
``unsupported by the admitted conditional route,'' leaving the original
question intact. Re-enabling native IDC requires a corrected or
otherwise justified path and an actual engine comparison against
independent collider references.

\hypertarget{a-front-door-reference-with-a-banking-interpretation}{%
\subsection{A front-door reference with a banking
interpretation}\label{a-front-door-reference-with-a-banking-interpretation}}

A second exact example supplies a positive identification reference. Let
\(A\) be assistance, \(M\) completion of the assisted process and \(Y\)
successful repayment. In the generated example an unobserved \(U\)
affects both assistance and repayment. The assumed graph has
\(A\to M\to Y\) and \(A\leftrightarrow Y\), with no direct effect of
assistance on repayment outside \(M\). This is a strong exclusion
assumption. In a real service workflow an agent's advice, fees or other
treatment channels could violate it, and observed associations alone do
not establish it.

Specify the finite law by

\[
 \begin{aligned}
 P(U=1)&=\tfrac12,\\
 P(A=1\mid U=u)&=\tfrac15+\tfrac35u,\\
 P(M=1\mid A=a)&=\tfrac14+\tfrac12a,\\
 P(Y=1\mid M=m,U=u)&=\tfrac1{10}+\tfrac12m+\tfrac15u.
 \end{aligned}
\]

Independent rational enumeration first constructs the observed joint
table, hiding \(U\), and evaluates

\[
 \sum_m p(m\mid a)\sum_{a'}E(Y\mid m,a')p(a').
\]

It separately evaluates the intervention law from the specified
mechanisms. Both routes produce \(13/40=.325\) under \(a=0\) and
\(23/40=.575\) under \(a=1\), for a difference of \(1/4\). The naive
observed difference is \(37/100\). The exact calculation therefore
distinguishes a successful symbolic formula from a shortcut that merely
compares observed groups. No fitted sample, randomized qualification
study or engine runtime is involved. \citep{exec13_y0_oracles}

If a proposal scenario assigns an assumed USD 1,000 benefit to each
additional successful repayment, the gross expected difference is USD
250 per account, before assistance costs and any other consequences.
That number is planning arithmetic under the declared law and value
assumption. It supplies neither observed bank ROI nor evidence that the
front-door graph describes a bank.

The inner treatment variable \(a'\) reveals an important serialization
requirement. It is a dummy summation coordinate, while the outer \(a\)
is the requested intervention value. A global string replacement could
capture the dummy variable and change the formula. The proposed result
format uses a typed expression tree with lexical bindings for
marginalized variables. It preserves probability factors, products, sums
and ratios, alpha-renames bound symbols when needed, and maps free
coordinates back to the original question. The inspected symbolic
classes expose those structures; rendered LaTeX or the library's
Python-like display string is suitable for presentation, not as an
executable interchange format. \citep{exec13_y0_dsl}

An independent negative reference is equally necessary. In a graph
allowing \(A\to Y\) and \(A\leftrightarrow Y\), take a fair hidden
\(U\), let \(A=U\), and let an independent bit \(E\) equal one with
probability \(.2\). The mechanisms \(Y=A\mathbin{\mathrm{xor}}E\) and
\(Y=U\mathbin{\mathrm{xor}}E\) give the same strictly positive observed
law of \((A,Y)\). Under intervention their effects are respectively
\(.6\) and zero. Both belong to the model allowing those edges; no
faithfulness restriction is imposed. A library cannot choose between the
effects from their identical observed distribution. A successful adapter
retains this nonidentification rather than translating every completed
request into a finite effect. \citep{exec13_y0_oracles}

\hypertarget{responses-deployment-and-the-next-prototype-decision}{%
\subsection{Responses, deployment and the next prototype
decision}\label{responses-deployment-and-the-next-prototype-decision}}

A symbolic identification result should contain the original query,
graph and context identities, engine revision, typed functional, support
conditions and source provenance. Its numeric estimate, standard error
and confidence region remain absent. When the unconditional ID path
reports nonidentification in an admitted model, the response should
identify that origin and retain its diagnostic without exaggerating its
verification. When the conditional joint route fails, the response is
unsupported. An invalid graph is rejected; a missing engine, internal
failure, timeout or cancellation is unavailable. These distinctions
determine whether the next action is to repair an input, add scientific
information, select a different method or provision a runtime.

Library, CLI and MCP surfaces should call the same core implementation
and return the same context and result identities. The symbolic service
requires graph and question artifacts, not bank account rows. A worker
can enforce wall-time, memory and output budgets while keeping engine
diagnostics separate from scientific conclusions. Cancellation must
preserve a completed result if publication has already occurred; an
interrupted partial expression cannot become an identified functional.
These are proposed implementation requirements, not performance
guarantees supplied by the inspected y0 API.

The pinned project metadata requires Python 3.11 or later and lists
twelve direct runtime dependencies, including NetworkX, SciPy, pandas,
scikit-learn and statsmodels. The build backend is \texttt{uv\_build} in
its declared version range. Its optional R extra uses \texttt{rpy2};
upstream test groups include additional packages that are not all
runtime requirements of the proposed pure-Python path. The inspected
tree has no root lockfile establishing resolved transitive versions. A
prototype admission therefore needs a frozen interpreter/platform and
dependency manifest, saved distribution hashes, actual build/import logs
and license notices. Source inspection alone cannot promise
compatibility with whatever environment happens to be installed.
\citep{exec13_y0_metadata,exec13_y0_tree}

The saved BSD three-clause license permits redistribution under its
notice, disclaimer and nonendorsement conditions. That supports a
possible integration route, while leaving dependency rights and actual
deployment behavior to be checked. It supplies no scientific warranty. A
port may become justified if a later measurement shows that the Python
environment or latency is a decisive constraint; there is no such
measured advantage for porting in the present evidence.
\citep{exec13_y0_license}

The prototype's first engine trials should compare the actual pinned
runtime with the independent front-door, bow and collider references,
including retained unsupported requests. Conversion checks should
exercise isolated variables, edge-type distinctions, coordinate
permutations, dummy-variable binding, zero conditional support and
distinct failures. A successful trial would support the named adapter in
its declared regime. It would not establish discovery accuracy, graph
truth, estimation calibration or institutional acceptance. The value of
this restricted design is that each later claim has a specific test and
a specific scientific object that must survive the interface. The next study tests that requirement inside graph modification, where a correctly converted isolated condition can still be lost; the later comparisons extend the same preservation question to custody, recovery and other engines.
\label{exec13:python-identification-end}

\label{exec13:interchange-end}


% SOURCE docs/release/2026-10-29/source/docs/writing-repair/2026-10-29/src/d8ecff6e9e360fae.tex CHAR 0
\section{An isolated condition must survive internal graph modification}
\label{exec13:isolated-condition}

Preserving the node set at the adapter boundary is necessary but does not
establish that every internal graph operation preserves it. The inspected y0
revision's \texttt{remove\_in\_edges} rebuilds a graph from the intervention
vertices and the endpoints of surviving edges. It does not explicitly seed
every original vertex. An isolated vertex outside the intervention can
therefore disappear. The ID recursion later requests ancestors of its outcome
set in this modified graph. These are source-level observations about the
same pinned revision, not results from an installed package trial
\citep{exec13_y0_graph,exec13_y0_id}.

Consider the supplied DAG with \(A\to Y\) and an isolated condition \(Z\).
The requested object is \(p(y\mid\operatorname{do}(a),z)\). The proposed
joint route begins with \(p(y,z\mid\operatorname{do}(a))\), retaining
\(Y,Z\) as its outcomes. All three vertices are ancestors of this outcome
set or members of it before graph modification, so the initial ancestor
restriction does not remove \(Z\). In the subsequent source path, removing
incoming edges to \(A\) reconstructs vertices \(A,Y\) from the intervention
and the edge \(A\to Y\). The ancestor request for \(Z\) then concerns a
vertex that is absent. An actual installed-engine run is required to record
its precise exception and response; the source trace already establishes
why checking only the initial conversion is insufficient
\citep{exec13_isolated_reference}.

This failure mechanism is distinct from a failure of causal identification.
For an exact reference, let \(A,Z\) be independent fair binary variables and
let an independent outcome mechanism satisfy
\[
 P(Y=1\mid\operatorname{do}(A=a))=\tfrac14+\tfrac12a.
\]
The eight observed cells have positive probability. Under the supplied
causal DAG and its independent mechanisms,
\[
 p(y,z\mid\operatorname{do}(a))=p(y\mid a)p(z),\qquad
 P(Y=1\mid\operatorname{do}(a),z)=\tfrac14+\tfrac12a
\]
on both supported values of \(z\). The effect is \(1/2\) within either
condition. Independent rational enumeration reproduces this result. These
are stipulated mechanisms, not inferred bank facts or fitted package
outputs \citep{exec13_isolated_reference}.

The admission decision is consequently more specific. The proposed joint
route needs tests for isolated conditions and outcomes at every internal
graph transformation, in addition to the front-door, bow and collider
references. A reviewed patch that retains all vertices would need an actual
engine comparison and regression checks before admission. An adapter may
instead implement an independently justified isolated-variable factorization
for the exact supported graph class, but must not discard a condition merely
because it appears inconvenient. A condition with an incident edge does not
qualify for the isolated-variable shortcut. Until a route passes these
checks, a missing-node exception is an unavailable computation with the
original question retained. It is neither a proof of nonidentification nor
a successful empty result. The public proposal can therefore recommend the
restricted integration while keeping its unresolved implementation obligation
visible.
\label{exec13:isolated-condition-end}


% SOURCE docs/release/2026-10-29/source/docs/writing-repair/2026-10-29/src/571c4aba469fd16e.tex CHAR 0
\hypertarget{numerical-equivalence-must-survive-composition}{%
\section{Numerical equivalence must survive
composition}\label{numerical-equivalence-must-survive-composition}}

Even an adapter that preserves the exact mathematical expression can
fail when independently acceptable numerical errors accumulate through a
workflow. Conditional identification provides a concrete example.
Suppose a returned functional is \(f=N/D\), with a strictly positive
denominator, and a downstream evaluator computes
\(\widetilde N,\widetilde D\). Assume certificates give

\[
 |N-\widetilde N|\leq\varepsilon_N,\qquad
 |D-\widetilde D|\leq\varepsilon_D,
 \qquad |N|\leq B,\quad D\geq d_0>\varepsilon_D.
\]

By writing the difference over a common denominator and using
\(|\widetilde D|\geq d_0-\varepsilon_D\), one obtains

\[
 \left|\frac{\widetilde N}{\widetilde D}-\frac ND\right|
 \leq\frac{\varepsilon_N}{d_0-\varepsilon_D}
 +\frac{B\varepsilon_D}{d_0(d_0-\varepsilon_D)}.
\]

Thus an absolute tolerance on probability calculations has no fixed
interpretation for the conditional answer as the denominator approaches
zero. The mathematical statement that a conditional probability exists
at a positive coordinate is weaker than a uniform numerical stability
claim. A service that promises a useful numeric answer needs either an
adequate support margin or a certified set that reflects the weak
denominator.

For a reproducible assumed example, let \(\widetilde N=.004\),
\(\widetilde D=.010\), with bounds \(\varepsilon_N=.001\) and
\(\varepsilon_D=.004\). Independent interval division gives

\[
 \frac ND\in\left[\frac{.003}{.014},\frac{.005}{.006}\right]
 =\left[\frac3{14},\frac56\right],
\]

while the point ratio is \(.4\). All numerator/denominator corners here
obey \(0\leq N\leq D\), so intersecting with the probability range does
not narrow the interval. A valid dependence relation between their
errors might narrow it, but it must be supplied and retained as joint
information. The broad result is produced by bounded deterministic error
alone; it is not a statistical confidence interval.
\citep{exec13_interchange_composition}

If an uncertainty region includes \(D=0\), the service should preserve
the unresolved support condition. Clipping the denominator to a positive
constant silently changes the functional and generally supplies no error
certificate. If only strictly positive but arbitrarily small
denominators are allowed, a ratio may be highly sensitive even though
the target is defined at every admitted point. The response should
distinguish a model-relative identification result, a supported numeric
coordinate and a finite-precision certificate. These objects answer
different questions and may become available at different times.

Composition also requires the correct order of nonlinear operations.
Suppose two equally weighted customer groups have success probabilities
\(1/5\) and \(1/2\). Their population probability is \(7/20\), with odds
\(7/13\). Averaging their odds instead gives \((1/4+1)/2=5/8\), which
transforms back to probability \(5/13\). Neither number arose from a
floating-point error. Averaging and the nonlinear odds map do not
commute. An interface that substitutes an average log-odds, odds or
probability because each is called a ``risk'' has changed the target. In
a causal contrast the error can affect both action-specific means before
they are differenced. \citep{exec13_interchange_composition}

These examples motivate a transformation graph accompanying the result.
Its nodes should identify the scientific object and coordinates, and its
arrows should identify an exact map, a conditional approximation or a
certified outer enclosure. A population aggregation followed by
valuation is one path; conversion to a subgroup score followed by
averaging is another. The workflow should be able to explain why the
selected path answers the declared question. A chain of individually
sensible package defaults is insufficient evidence that the composed
path does.

Cross-language comparison can test these properties without pretending
to qualify every estimator. For an exact finite reference table, an
admitted identification expression should agree with a separately
computed intervention distribution; variable renaming and category-order
changes should preserve that answer after inverse mapping. A sign
reversal should reverse an effect and reorder interval endpoints. A
probability marginal should normalize on its declared support. These are
necessary checks on the concrete path, not sufficient evidence of
universal scientific correctness. The front-door and collider examples
provide particularly useful references because a superficially plausible
formula can return a different probability.

The same testing discipline should retain expected refusals. An
unsupported graph class must not become an empty graph. A normalization
failure must not become an estimated zero effect. A numerical timeout
must not be cached as a scientific negative answer for an unrelated
context. Matching those dispositions across library, CLI and MCP is part
of semantic equivalence. The preserved local transport trial
demonstrates selected comparisons for the existing bounded operations;
the proposed y0 path still requires its own actual engine and conversion
evidence. \citep{exec13_interchange_runtime}

For investment planning, the practical question is whether the
information and computation are adequate for the decision. A numerically
precise conditional effect on a vanishingly small subgroup may be
irrelevant to the bank's target population. A broad but valid effect set
may establish that a costly expansion cannot be justified yet. A
slightly more expensive evaluator may be valuable if certified
containment resolves a decision that uncontrolled approximations cannot.
The proposed product should expose these consequences through the
scientific result, its restrictions and its value calculation, rather
than hiding them behind a generic success flag.
\label{exec13:interchange-complete-end}


% SOURCE docs/release/2026-10-29/source/docs/writing-repair/2026-10-29/src/df465d8d642959cb.tex CHAR 0
\section{Preserving graph classes, coordinates and set geometry}
\label{exec16:topology}

A bank can receive a syntactically valid result for a different decision from
the one it submitted. A graph may retain its adjacency while losing an
unresolved endpoint. Two effects measured in dollars may exchange positions.
A disconnected confidence set may become a convenient interval that includes
values the calculation excluded. These changes matter before an estimator's
speed or predictive accuracy is compared. The proposed integration layer must
therefore preserve the object on which the decision depends: the scientific
context, the graph or graph class, the named coordinates, and the meaning and
geometry of the returned set. A transport test should deliberately construct
examples in which losing one of these features changes an answer.

The distinction between a declaration and its scientific validation is
especially useful here. The local reference service already rejects directed
cycles in a declared DAG, ambiguous DAG endpoints, changed context hashes,
inconsistent set topologies, and several unsupported identification requests.
Its graph-review operation explicitly reviews syntax. Direct calls also accept
an ADMG declaration containing a circle endpoint or a directed cycle, and a
CPDAG declaration with a bidirected edge, while returning the limitation that
full graph-class membership has not been tested. These are observed boundaries
of the selected operation, not evidence that those objects belong to the named
classes. The separate problem-declaration reviewer rejects the tested ADMG
circle and performs additional mark, duplication, cycle and event-order checks.
It still delegates full CPDAG and PAG recognition to a formal engine. An adapter
must retain the review level as well as the graph label; neither a successful
request nor the word \texttt{computed} is a graph-theoretic certificate.
\citep{exec16_topology_runtime}

\subsection{A class of graphs is not their union}

Consider the three-node chain with variables \(A,M,Y\), with no adjacency
between \(A\) and \(Y\). The three orientations
\[
 A\longrightarrow M\longrightarrow Y,\qquad
 A\longleftarrow M\longrightarrow Y,\qquad
 A\longleftarrow M\longleftarrow Y
\]
have no unshielded collider. For a positive observed law satisfying
\(A\mathrel{\perp\!\!\!\perp}Y\mid M\), their observational factorizations agree:
\begin{align*}
 p(a)p(m\mid a)p(y\mid m)
 &=p(m)p(a\mid m)p(y\mid m)\\
 &=p(y)p(m\mid y)p(a\mid m).
\end{align*}
This equality is an application of conditional independence and Bayes' rule.
It does not declare the corresponding interventions equivalent. An exact
enumeration of all directed simple graphs on these three labelled vertices
finds 25 DAGs and 11 groups defined by their skeletons and unshielded collider
patterns. The chain group contains these three members. Taking the union of
their directed edges introduces arrows in both directions on the adjacent
pairs, producing cycles that occur in none of the members. Storing that union
as a DAG thus destroys the object being represented.
\citep{exec16_topology_exact}

A set representation should preserve member identities or a precisely defined
equivalence-class representation. An endpoint mark must retain its declared
meaning; adjacency alone cannot preserve the distinction between a tail,
arrowhead and circle. Additional background restrictions also need their own
identity, because they can reduce the admissible orientations. A declared
member handle is only a reference. In an actual syntax-review probe, an absent
member with a well-formed hash and handle is counted and returned without
resolving its bytes. The proposed class-review operation must separately
resolve every required member, verify its content and schema, check its
variable map, and record the exact recognition domain. Missing members require
an unavailable review, not silent deletion from the class.
\citep{exec16_topology_runtime}

The effect consequence can be seen without a large simulation. Let independent
zero-mean Gaussian errors generate either
\[
 A=\varepsilon_A,\quad Y=\tfrac12 A+\varepsilon_Y,
 \quad (\operatorname{Var}\varepsilon_A,\operatorname{Var}\varepsilon_Y)
       =(1,3/4),
\]
or
\[
 Y=\widetilde\varepsilon_Y,\quad A=\tfrac12 Y+\widetilde\varepsilon_A,
 \quad (\operatorname{Var}\widetilde\varepsilon_Y,
        \operatorname{Var}\widetilde\varepsilon_A)=(1,3/4).
\]
Both observed vectors \((A,Y)\) have covariance
\(\left(\begin{smallmatrix}1&1/2\\1/2&1\end{smallmatrix}\right)\), hence the
same Gaussian law. Increasing \(A\) by one through intervention changes the
mean of \(Y\) by \(1/2\) in the first model and zero in the second. At a
stipulated USD 20 per outcome unit and USD 8 action cost, the net values are
USD 2 and USD \(-8\). Choosing one orientation because a receiving library
requires a DAG would choose between these decisions without new evidence.
The appropriate proposal is to retain both effects or obtain a defensible
restriction, intervention or additional measurement that distinguishes them.
\citep{exec16_topology_exact}

The supplied-DAG back-door operation in these preserved trials does not implement that general class calculation. The later restricted Gaussian research discussed below is separate evidence; it does not widen the tested operation's contract. The actual probes return \texttt{unsupported} for PAG,
CPDAG, ADMG and DAG-set requests. A supplied DAG with an unblocked latent
common cause also returns an unsupported back-door attempt, explicitly
distinguished from a general nonidentification result. A future identification
engine must declare which graph family and criterion it handles. For a
retained model collection \(\mathcal G\), a proposed aggregate effect region is
\(\bigcup_{G\in\mathcal G}\Theta_G\), where \(\Theta_G\) includes all effects
compatible with the stated observational law and assumptions within member
\(G\). This expression requires each constituent calculation to have the
claimed meaning. An uncomputed constituent cannot be dropped or converted to
the empty set. Neither a graph-class representation nor a conditional
identification formula supplies sampling calibration by itself.
\citep{exec16_topology_runtime}

\subsection{Coordinate identity is stronger than a unit label}

Suppose retail and corporate effects are USD 40 and USD \(-10\) per eligible
account. The requested population weights are \(4/5\) and \(1/5\), and the
action costs USD 5. The net effect in the declared order is
\[
 \tfrac45(40)+\tfrac15(-10)-5=25.
\]
If a response uses corporate--retail order but a consumer applies the weights
by position, the result becomes
\[
 \tfrac45(-10)+\tfrac15(40)-5=-5.
\]
Both vectors have length two and both coordinates use USD. The existing schema
and semantic checks accept an authored response whose generic payload declares
the reversed order, while the requested population text declares retail first.
They check units and coordinate counts, but do not bind an explicit coordinate
identity vector between these locations. This observation concerns an admitted
ordering mismatch; it does not claim that any schema can prove that submitted
numbers are scientifically correct.
\citep{exec16_topology_runtime,exec16_topology_exact}

The proposed contract should give each coordinate a stable identifier binding
its outcome, intervention contrast, population, horizon, statistic and units.
Results should pair these identifiers with their values and set components.
A consumer either matches the requested order or applies an explicit declared
permutation. If \(P\) is its permutation matrix, the transformation must act on
the estimate, covariance, joint-region constraints and financial weights
together. For a general fixed affine map \(\phi=L\theta+c\), the familiar
relations
\[
 \widehat\phi=L\widehat\theta+c,\qquad
 \widehat\Sigma_\phi=L\widehat\Sigma L^\top
\]
do not license a change in the causal target. The map needs its own input and
output identities, and an estimated map needs its own uncertainty treatment.
The earlier population and covariance examples establish why this binding
must include the weighting measure and joint errors.
\citep{exec13_interchange_examples}

Stable identifiers make an accidental reordering detectable when values retain
their identities. They cannot detect a producer that attaches a wrong value to
the correct identifier. Native adapter acceptance therefore also needs
asymmetric sentinel examples, deliberate node and column permutations, and
independent numerical references. Equal effects or identical units are poor
sentinels because they can conceal the error. The proposed retail example is
small enough to run at every language boundary before any financial trial.

\subsection{Set images and joint restrictions}

Unknown, empty, whole-line and disconnected sets have different meanings.
The current semantic representation distinguishes an uncomputed unknown set
with \texttt{intervals:null}, an empty set with \texttt{intervals:[]}, and the
whole line with a single pair of unbounded endpoints. The probes preserve
these states, finite unions and singleton intervals, and reject inverted,
overlapping or inconsistently labelled pieces. A refusal with no numerical
result remains separate from all of them. In particular, an unsupported
identification algorithm does not establish an empty identified set.
\citep{exec16_topology_runtime}

For a scalar map \(g(x)=1-2x\), consider
\[
 S=(-\infty,-2]\cup[2,\infty).
\]
Because the slope is negative, the image reverses both endpoint order and
component order:
\[
 g(S)=(-\infty,-3]\cup[5,\infty).
\]
For the finite union \([-3,-2]\cup[4,5]\), the image is
\([-9,-7]\cup[5,7]\). Replacing it by its convex hull \([-9,7]\) introduces
the entire gap \((-7,5)\). The existing scalar interval-union arithmetic
reproduces these images, including sign reversal and unbounded endpoints.
An outer approximation can be appropriate if it is explicitly labelled and
its consequences are understood. It cannot be described as the exact image
when it fills excluded values.
\citep{exec16_topology_exact}

Coordinate projections lose a different kind of information. Let a stipulated
joint region for two dollar effects be
\[
 C=\{(x,10-x):0\leq x\leq10\}.
\]
Each projection is \([0,10]\), but every point in \(C\) has sum ten. With
USD 12 cost the exact joint net-value image is the singleton \(\{-2\}\).
Applying interval-sum arithmetic to the Cartesian product of the projections
instead yields \([-12,8]\). The arithmetic correctly computes the image of
the larger rectangle. It has not recovered the relation shared by the two
coordinates. A consumer seeking a robust acceptance decision would reject
this action under both sets, but a possibility-of-profit screen admits it
under the rectangle and rejects it under the actual joint region. Numerical
agreement on coordinate endpoints therefore does not establish agreement on
the proposed decision.
\citep{exec16_topology_exact}

The contract should identify whether a set is a marginal projection, a product
outer bound, or an encoded joint region, and preserve any constraints needed
by the downstream transformation. Sampling meaning is another independent
field. Two pointwise 95\% sets provide a joint coverage lower bound of only
90\% by the union bound; under independent exact coverage events their joint
coverage is \(0.95^2=0.9025\). A declaration of simultaneous coverage requires
a construction that supports it. The local validator requires construction
evidence for a joint-confidence label, but cannot certify the validity of
that evidence. The selected three-variable independent Gaussian discovery
construction also returns conservative outer bounds under its restricted
linear, causally sufficient model class. Its existence does not qualify
arbitrary temporal PAGs or a general graph-to-effect region.
\citep{exec16_topology_runtime,exec16_topology_exact,exec16_topology_source}

\subsection{A lossless-looking JSON round trip}

An actual Python--Node--Python parse-and-serialize probe preserves the tested
endpoint marks, graph-member hash string, unknown and empty states, whole-line
and disconnected sets, and a large identifier stored as a string. This is
evidence about this generic codec path; no R, Java or C++ causal adapter was
executed. The same path changes an integer estimate and singleton endpoint
from \(2^{53}+1=9007199254740993\) to
\(2^{53}=9007199254740992\). Both the original and changed response pass the
present schema and semantic checks. The resulting singleton excludes the
original value despite preserving the apparent topology and context hash.
\citep{exec16_topology_codec}

The example is deliberately chosen at an exactness boundary; it does not
claim that ordinary bank effects have that magnitude. Large integer monetary
counts, identifiers and scaled fixed-point quantities can nevertheless cross
language boundaries. A prototype should either reject integers outside a
declared exact numeric domain or introduce a typed decimal/integer string
representation with a scale and rounding rule. The present numeric endpoint
schema would need an explicit successor to admit that representation.
Infinities should retain their structured endpoint meaning, rather than being
silently replaced by large finite sentinels. For approximate endpoints, the
contract must state error bounds and, where set inclusion is claimed, use
justified outward rounding.

Context hashing exposes a complementary failure. Python's current sorted-key
JSON digest distinguishes a horizon value serialized as \texttt{12.0} from
one serialized as \texttt{12}. The Node round trip changes the former to the
latter, after which the reference validator rejects the unchanged response
hash. Here the number has the same mathematical value, but the byte convention
differs. This conservative rejection is preferable to accepting an actually
changed question; a shared canonical encoding is still necessary for reliable
interoperability. The proposed encoding must specify number forms, strings,
key ordering and exact type distinctions. Alternatively, an adapter can retain
the original canonical context bytes as an opaque object and hash those bytes.
Simply recomputing a digest after an arbitrary transformation would hide the
very changes that context binding is intended to detect.
\citep{exec16_topology_codec}

\subsection{A bounded prototype acceptance contract}

The first implementation comparison should use generated objects with exact
answers and documented public examples. It should admit a narrow graph family,
resolve and validate all members, and compare graph representations under
declared node renaming. The three-node enumeration supplies a small exhaustive
reference; additional tests must cover each supported family, background
restriction and latent/temporal convention. Recognition, observational
equivalence and intervention identification should have separate result types.
Unsupported families must remain visible in the request denominator.

At the numerical boundary, the same asymmetric coordinates and interval unions
should pass through each proposed native adapter, library, CLI and MCP route.
Acceptance requires preservation of identities and topology or an explicitly
bounded approximation, agreement with independent affine and joint-region
references, and rejection of undeclared rounding or reindexing. Deliberately
missing members, changed hashes and unknown constituent effects should produce
typed unavailable or unsupported results. Deadline and resource failures must
not publish a partially assembled graph class as complete. This comparison
has a known engineering route, but the generic JSON probe alone does not
establish any native adapter's acceptance.
\citep{exec16_topology_contract}

The harder scientific question remains separate: a retained observational
class may imply several intervention effects, and a learned class need not
have calibrated retention probability. Numerical preservation can prevent
additional distortion; it cannot supply missing causal information or a new
uncertainty theorem. The proposal should select a model and sampling regime,
justify its identification and joint construction, and use generated truth to
study that regime before proposing expensive bank experiments. Later bank
owners must supply variable meanings, feasible interventions, clocks,
background restrictions and decision tolerances. Platform and consumer owners
must accept their actual encodings and failure behavior. These are subsequent
acceptance inputs, while the public/generated prototype and the present
proposal comparisons can proceed now.
\citep{exec16_topology_contract}
\label{exec16:topology-end}


% SOURCE docs/release/2026-10-29/source/docs/writing-repair/2026-10-29/src/05e0c3a12e3de929.tex CHAR 0
\hypertarget{data-custody-preserves-the-meaning-of-a-causal-claim}{%
\section{Data custody preserves the meaning of a causal
claim}\label{data-custody-preserves-the-meaning-of-a-causal-claim}}

\label{exec14:artifact-custody}

\hypertarget{the-object-to-preserve-is-a-decision-with-its-information}{%
\subsection{The object to preserve is a decision with its
information}\label{the-object-to-preserve-is-a-decision-with-its-information}}

A servicing analysis can be arithmetically reproducible and still answer
the wrong question. Suppose an analyst computes the effect of an offer
on the next twelve months of account value among accounts eligible at
the decision date. A second workflow receives the same numerical moments
but reports a twenty-four-month effect for all customers. No file
corruption is needed for that mistake. Population, outcome horizon,
intervention, information clock, units and sampling law belong to the
scientific object being transferred. They must travel with the data,
graph and result. A consumer that cannot establish this agreement should
return an unavailable result or request an explicit new analysis.

Let \(C\) denote the declared scientific context, including those fields
and the graph with its assumption statuses. Let \(S\) be the local
deterministic JSON serialization, \(H\) the SHA-256 digest, and
\(B=S(V)\) the stored bytes for a JSON value \(V\). The implemented
local store uses a descriptor

\[
D=(H(B),\;\text{media type},\;\text{schema},\;\text{role},\;|B|,\;\text{custody}),
\]

and derives a handle from \(H(S(D))\), preceded by a fixed letter. A
scientific input also carries \(H(S(C))\) inside its payload. These are
separate bindings: the first identifies bytes and their declared
interpretation, while the second ties a payload to a particular
scientific request. The store's serialization sorts object keys and
excludes nonfinite JSON numbers. It is an inspected Python serialization
rule, not an established canonicalization standard for every language or
a proof that distinct textual descriptions express the same causal
question.~\citep{exec14_custody_implementation}

For instance, a moment summary contains both its numbers and the context
digest. Changing the horizon from twelve to twenty-four months changes
the context digest even when all moments are unchanged. The service
checks the expected artifact role and schema, reads and verifies the
referenced bytes, and compares the embedded context digest. In an actual
local trial with a generated quadratic sensitivity input, the
twelve-month request computed an empty sublevel set; the otherwise
identical twenty-four-month request raised a context error. Through the
subprocess runtime the latter became a failed, unavailable response with
no numerical result. This distinction matters: an empty set is a
mathematical outcome for a supplied problem; an unavailable result says
that the requested problem was not computed. Conflating them would turn
a custody failure into apparent scientific evidence.
\citep{exec14_custody_observations}

Hash agreement does not establish that the population declaration is
true, that timestamps reflect actual availability, or that a supplied
graph is causal. The declaration and the authority behind it require
different evidence. A data owner can provide a signed account of
extraction and eligibility, but the signature would establish who made
that account, not the validity of ignorability, exclusion or a
measurement model. The proposed product should preserve these
distinctions from ingestion through the reported monetary consequence.

\hypertarget{what-the-local-implementation-establishes}{%
\subsection{What the local implementation
establishes}\label{what-the-local-implementation-establishes}}

The reference store is restricted to nonsymlink runtime paths inside
this repository. It stores finite JSON values up to eight MiB and
publishes each content or metadata file atomically, with a directory
flush. Reading checks agreement between the supplied reference and
recorded metadata, file length and the byte digest. The operation
dispatcher adds role, schema and declared scientific-context checks. The
runtime verifies input references before execution and before
considering an idempotent replay. \citep{exec14_custody_implementation}

Fifteen generated local cases exercise this boundary. Repeated insertion
under the same role returned the same reference; changing only the role
retained the content digest and changed the handle. Changed bytes,
changed metadata, an incompatible horizon, a wrong role or schema, an
insufficient read limit, a nonfinite number, and selected symlink paths
were refused. The two existing service/runtime test modules also passed
all sixteen tests in this environment. These observations support
selected implementation claims; they do not estimate an operational
error rate, prove correctness against every possible input, or validate
bank records. \citep{exec14_custody_observations}

Replay deserves its own scientific interpretation. A new request
correlation identifier may retrieve the same result, but reusing the
idempotency key for a changed horizon returned a conflict. After the
referenced input was altered, replay returned an artifact-access failure
with no result. Its key binds the request digest, selected
implementation and schema hashes, Python version and declared dependency
versions. A current source digest prevents an ordinary unchanged-key
replay across a detected implementation change; it is not a complete
software supply-chain attestation. The local operator still controls
these files and the environment.
\citep{exec14_custody_implementation,exec14_custody_observations}

The report path has an equally narrow role. Its summaries explicitly
carry a schema-only verification label. Preserving a canonical response
in a report does not promote execution evidence into accepted causal
inference. For example, an assumption-conditional identification
expression must retain that qualification even if its artifact passes
every byte check. Similarly, a refused or unsupported operation must not
disappear from a report that counts only successful tasks. The preserved
runtime tests establish selected report and cancellation behavior;
institutional report approval remains an additional
decision.~\citep{exec14_custody_implementation,exec14_custody_observations}

\hypertarget{two-observed-boundaries-and-the-repairs-they-require}{%
\subsection{Two observed boundaries and the repairs they
require}\label{two-observed-boundaries-and-the-repairs-they-require}}

The first boundary concerns the relation between a handle and its
descriptor. The writer derives the handle from the descriptor, but the
current reader does not independently recompute that derivation. A
controlled local trial changed a generated balance from 100 to 900,
rewrote the stored descriptor with the new byte digest and length, and
supplied the correspondingly changed reference under the original
handle. The old trusted reference was refused. The new supplied
reference was accepted, although recomputing the descriptor digest
produced a different handle. The test required write access to content
and metadata and control of the supplied reference. It demonstrates a
limit of the single-operator model; it does not establish an exposed
remote attack path.~\citep{exec14_custody_observations}

There is a direct local repair: validate the complete descriptor schema
and recompute its handle before returning content or accepting stored
metadata. Require the derived handle to equal both the requested
identifier and the identifier in the stored reference. Preserve the
original implementation as the evidence baseline, and qualify a
successor against the old trusted reference, a coherently rewritten
reference, a mismatched metadata identity, missing descriptor fields and
extra fields. This is a proposed repair, not an implemented feature. It
also does not solve unrestricted substitution: a writer who can replace
the trusted request can supply an entirely new, internally consistent
artifact and handle. Preventing that substitution requires an
independently trusted manifest or authenticated authorization binding
the requester, purpose and allowed inputs. \citep{exec14_custody_design}

The second boundary concerns availability. The reader opens a path
before checking that its file descriptor denotes a bounded regular file.
The open uses a flag rejecting a final symlink, but lacks nonblocking
mode. In a disposable subprocess, opening a named pipe without a writer
did not return within a one-second deadline after the child signalled
that it was about to read. The subprocess was terminated. The observed
block occurred before the regular-file test could reject the pipe.
Because the runtime performs its initial reference verification in the
parent, a deadline on the later worker does not by itself bound that
initial access. This is an actual local availability finding, with no
claim about deployment reachability or a measured service-level failure
rate. \citep{exec14_custody_implementation,exec14_custody_observations}

A candidate repair opens candidate paths without following a final
symlink and in nonblocking mode, immediately checks the returned
descriptor, and accepts only bounded regular files. The acceptance cases
must include pipes with and without writers, directories, symlinks and
oversized or changing files; each rejection must occur before any causal
worker starts. Directory file descriptors and controlled directory
ownership should also be reviewed if the threat model includes
concurrent path replacement. A prior path resolution check alone does
not establish race-free access. Nonblocking open is a concrete first
repair for the demonstrated pipe, not a promise to bound all filesystem
delays or solve every race. A production design still needs an access
deadline and an operating-system-specific review.
\citep{exec14_custody_design}

These findings do not invalidate the observed ordinary mutation
refusals. They define the assumptions under which those refusals are
useful. For a single researcher working with generated inputs, agreement
with a trusted reference can detect accidental changes. For multiple
operators, automated agents or a bank repository, the trust anchor and
authorization decision must live beyond a request's freely supplied
custody string.

\hypertarget{custody-across-languages-repositories-and-agent-tools}{%
\subsection{Custody across languages, repositories and agent
tools}\label{custody-across-languages-repositories-and-agent-tools}}

The transfer contract should distinguish a logical artifact from a path
on one machine. A Python caller, an R adapter and a Java discovery
engine may share a causal task without sharing a filesystem. The
prospective adapter resolves an allowed handle within its own custody
boundary, checks the manifest and exact schema, and obtains the material
needed for its declared operation. It preserves graph endpoints,
effect-set topology, units, population, clock and uncertainty semantics
in the returned object. Copying a path into a CLI flag or an MCP request
does not establish that agreement. No new sibling-repository execution
is claimed by this design. \citep{exec14_custody_design}

Two prospective arrangements have different costs. A controlled local
service may expose a small allowlist of typed handles and operations to
trusted callers. This keeps initial experiments inexpensive but depends
on the service's operator and host boundary. A broader artifact service
would bind authenticated principals and permitted purposes to manifests,
retain authorized provenance, and enforce current access and revocation
decisions. That arrangement adds operational complexity and requires an
owner for identity, keys, access policies and lifecycle events. The
product proposal can specify and compare them before claiming either is
an accepted bank service.

A context digest provides a strict equality check, which can be
desirable at a computational boundary but does not resolve semantic
equivalence. Equivalent population wording or a different serialization
of the same declared numerical quantity may fail equality; different
underlying data can also carry an identical but false declaration. A
future interoperable context schema therefore needs a versioned
representation, explicit units and clock conventions, and reviewed
migrations. If an adapter changes a horizon, aggregation or currency
conversion, it should produce a new artifact with the transformation and
input identities recorded. It should not overwrite the meaning attached
to an existing result.

Retention creates another distinction between repeatability and
permission. An immutable result can remain scientifically reproducible
while its source records become unavailable for a later request. The
service should check current access policy before replay, distinguish
revoked access from missing bytes, and retain only the audit material
its owner permits. An invalidated source should have an explicit status
and downstream impact record. Deletion, backups, redaction, derived
artifacts and legally permitted audit retention require owner decisions.
The current local store has no qualified revocation or institutional
retention service. Designing one requires more than continuing to
compare hashes.~\citep{exec14_custody_design}

\hypertarget{an-affordable-qualification-path}{%
\subsection{An affordable qualification
path}\label{an-affordable-qualification-path}}

The next experiment can use generated accounts, public documents and
deliberately altered copies. It should first qualify the two
demonstrated local repairs without changing the original evidence. Then
it can exercise an adapter against distinct populations, horizons,
endpoint types and information clocks, including refusals and missing
inputs. Public data can test acquisition and transformation lineage, but
a public label does not automatically establish redistribution rights or
causal ground truth. Each chosen dataset needs its own provenance and
permitted-use review.

The prospective acceptance cases divide into concrete responsibilities:

\begin{longtable}[]{@{}lll@{}}
\toprule
\begin{minipage}[b]{0.30\columnwidth}\raggedright
Requirement\strut
\end{minipage} & \begin{minipage}[b]{0.30\columnwidth}\raggedright
Local or public-evidence next step\strut
\end{minipage} & \begin{minipage}[b]{0.30\columnwidth}\raggedright
Additional owner input\strut
\end{minipage}\tabularnewline
\midrule
\endhead
\begin{minipage}[t]{0.30\columnwidth}\raggedright
Byte, descriptor and handle consistency\strut
\end{minipage} & \begin{minipage}[t]{0.30\columnwidth}\raggedright
Qualify a successor reader against ordinary and coherent replacement
cases\strut
\end{minipage} & \begin{minipage}[t]{0.30\columnwidth}\raggedright
Trusted manifest location and who may publish it\strut
\end{minipage}\tabularnewline
\begin{minipage}[t]{0.30\columnwidth}\raggedright
Bounded input access\strut
\end{minipage} & \begin{minipage}[t]{0.30\columnwidth}\raggedright
Reject unexpected file types before worker launch; test deadlines and
controlled path replacement\strut
\end{minipage} & \begin{minipage}[t]{0.30\columnwidth}\raggedright
Target filesystem, latency limits and operating model\strut
\end{minipage}\tabularnewline
\begin{minipage}[t]{0.30\columnwidth}\raggedright
Scientific meaning\strut
\end{minipage} & \begin{minipage}[t]{0.30\columnwidth}\raggedright
Preserve role, schema, context, endpoint and uncertainty contracts in
adapter fixtures\strut
\end{minipage} & \begin{minipage}[t]{0.30\columnwidth}\raggedright
Population, time, assignment and measurement declarations with
accountable provenance\strut
\end{minipage}\tabularnewline
\begin{minipage}[t]{0.30\columnwidth}\raggedright
Authorization and agent scope\strut
\end{minipage} & \begin{minipage}[t]{0.30\columnwidth}\raggedright
Model permitted and forbidden handles, purposes and operations using
generated principals\strut
\end{minipage} & \begin{minipage}[t]{0.30\columnwidth}\raggedright
Identity authority, policy owner, key custody and access rules\strut
\end{minipage}\tabularnewline
\begin{minipage}[t]{0.30\columnwidth}\raggedright
Retention and revocation\strut
\end{minipage} & \begin{minipage}[t]{0.30\columnwidth}\raggedright
Exercise expired and revoked input scenarios without silently replaying
results\strut
\end{minipage} & \begin{minipage}[t]{0.30\columnwidth}\raggedright
Retention, deletion, backups and permitted audit evidence\strut
\end{minipage}\tabularnewline
\begin{minipage}[t]{0.30\columnwidth}\raggedright
Cross-repository use\strut
\end{minipage} & \begin{minipage}[t]{0.30\columnwidth}\raggedright
Compare actual local behavior with the declared adapter contract\strut
\end{minipage} & \begin{minipage}[t]{0.30\columnwidth}\raggedright
Named consumer, supported environment and acceptance decision\strut
\end{minipage}\tabularnewline
\bottomrule
\end{longtable}

The bank need not fund an authentic institutional trial to establish the
value of these checks. A prototype can first show that a wrong horizon
is refused, an unavailable result stays unavailable, and failed access
cannot masquerade as an empty effect set. Its success criterion is
preserved meaning and a documented disposition for every request. Actual
bank provenance, access authority, operating performance and adoption
remain later claims requiring their own
evidence.~\citep{exec14_custody_design}

\label{exec14:artifact-custody-end}


% SOURCE docs/release/2026-10-29/source/docs/writing-repair/2026-10-29/src/ddab88fda0ff15a7.tex CHAR 0
\section{Recovery that preserves a causal request}
\label{exec14:failure-trials}

A recoverable service must preserve the scientific request as well as the
computation. Suppose a servicing analyst requests an effect in dollars over
twelve months and the calculation is interrupted. A later response for euros,
another horizon or a revised graph is not a continuation of that request.
Conversely, the absence of a response says nothing about whether the effect is
zero, unidentifiable or economically unattractive. Our proposed integration
therefore maintains separate identities for the scientific context, execution
environment, input artifacts, completed work and published result.
\citep{exec14_failure_design}

The local reference runtime already binds requests to context and artifact
identities, and binds replay to selected code, schemas, Python and dependency
versions. Its report jobs preserve a prefix of completed artifact identifiers
and corresponding summaries. Reports describe supplied results as declarations
checked against a schema; they do not confer independent causal review. These
are useful mechanics for an early prototype, with a deliberately narrow
single-operator scope. They neither authenticate the supplied business premises
nor qualify the system for bank operation. \citep{exec14_failure_implementation}

\subsection{Semantic failure and operational failure require different responses}

Generated local requests illustrate the distinctions. A valid quadratic
sensitivity request returned a computed empty set. Changing its declared units
from USD to EUR while retaining the original input artifact produced a failed,
unavailable response with no mathematical result. An unsupported discovery
method returned a completed invocation with the scientific disposition
unsupported and no result. A one-byte output limit produced an output-budget
failure. These observations demonstrate distinct response meanings for selected
requests; an empty mathematical set must not be used as a generic representation
of missing computation. \citep{exec14_failure_observations}

The review also exposed a boundary that the proposed interface must repair.
The library raises a validation exception for an unsupported request-schema
version. It also raises a JSON decoding exception when a previously saved replay
file is manually truncated to an incomplete JSON object. The latter case used
an unchanged, valid request; the fault was in local saved state. A follow-up
through the actual CLI found that both conditions return exit code two and
the error code invalid\_request. The CLI catches the exception, but assigns
the storage fault to the wrong operational category. A caller could repeatedly
rewrite a scientifically valid request while the corrupted replay state remains
unchanged. \citep{exec14_failure_observations}

This finding does not establish that the runtime's atomic writer naturally emits
truncated files. The truncation was a deliberate local fault injection. It does
establish the observed reader behavior for those bytes, and a concrete recovery
requirement. The proposed response should classify invalid caller syntax
separately from corrupt or unavailable saved state, preserve the offending bytes
and their identity for diagnosis, and expose whether retry can succeed. It should
not silently delete the cache and rerun a scientific operation whose original
inputs, environment, randomization or authorization may no longer be available.
A rebuild must first revalidate those identities and declare whether it resumes
the old request or creates a new attempt. \citep{exec14_failure_design}

\subsection{Three interruption windows}

The content-addressed artifact store publishes bytes and metadata through
separate atomic writes. Individual write atomicity does not make the pair one
transaction. In a generated case we paused immediately after publishing the
content file, before publishing its descriptor, and killed that child process.
The resulting content-only artifact was absent from the reference listing and
could not be retrieved by descriptor lookup. Retrying the same value, schema
and role completed the missing metadata under the same handle and returned the
original value. This supports one idempotent repair route. A different value
under that handle would require a different decision; this observation does not
authorize silently reusing a stale descriptor. \citep{exec14_failure_observations}

The second window concerns work already completed but not yet delivered.
A report over three generated canonical responses was interrupted after its
first saved summary. An unmodified runtime then resumed the job. The final
ordered identifiers exactly matched the original three requested artifacts,
and the report contained one summary per identifier. Every summary retained
the qualification that checking was schema-only, without causal or external
review. Preserving that qualification across recovery is as necessary as
avoiding duplicated work: a restart must not upgrade a conditional result into
an accepted finding. \citep{exec14_failure_observations}

The third window occurs after the report artifact has been stored but before the
job's terminal state points to it. We paused at that exact boundary and killed
the child. Saved job state still reported running with no result reference,
although the report artifact already existed. Recovery completed the job and
returned the same report handle and bytes. Thus a durable content-addressed
artifact can support idempotent publication even when acknowledgement is lost.
The demonstrated behavior concerns local deterministic summaries. A future
external side effect, such as submitting a bank instruction, would need its own
transaction and reconciliation protocol; repeating a local report write supplies
no evidence about that different operation. \citep{exec14_failure_observations,exec14_failure_design}

These three trials used explicit test-process hooks at named publication
boundaries. The library source files were unchanged. The test child executed
the actual store or report-worker code up to the boundary, and the parent killed
only its own child. Subsequent recovery used the ordinary runtime. Initial
submission scheduling was deliberately deferred so the selected boundary could
be reached deterministically; therefore the production spawn counter excludes
that injected first worker. We do not treat the displayed counter as a count
of every process launched by the harness. These cases give reproducible
mechanism evidence, not outage frequencies or coverage of every operating-system
failure ordering. \citep{exec14_failure_observations}

\subsection{A checkpoint is evidence to validate, not an instruction to trust}

In another generated case, the first saved summary was changed while retaining
its old checkpoint digest. Resume failed with a checkpoint-content mismatch,
kept the result reference absent and preserved the failed state. This is the
appropriate local refusal for that mutation. The implementation also checks
that completed identifiers form a prefix of the frozen request, that summary
and identifier counts agree, and that execution identity remains compatible.
These checks bind a particular continuation. They do not authenticate a writer
who can coherently replace every relevant descriptor and trust anchor; the
separate custody discussion retains that limitation. \citep{exec14_failure_implementation,exec14_failure_observations}

The proposed state transition can be expressed without assuming success:
\[
 \text{queued}\longrightarrow\text{running}
 \longrightarrow\{\text{completed},\text{failed},\text{cancelled}\}.
\]
Resumption of unfinished work is allowed only after the input, execution and
checkpoint conditions hold. Completion additionally requires a readable result
whose content and meaning match the declared request. Cancellation records a
terminal decision about future work; it is not a claim that previously published
artifacts or external effects have been rolled back. A nonterminal job may
already have a durable artifact, as the third interruption case demonstrates.
The authoritative job state and the artifact's independent identity must remain
distinguishable. \citep{exec14_failure_design}

For the next prototype, the known repair is to centralize error normalization
for library, CLI and MCP callers, while preserving the difference between
request defects, state corruption, resource exhaustion and scientific refusal.
Acceptance cases should include malformed and missing replay records, wrong
context or schema, unreadable referenced artifacts, changed code identity,
partial metadata, corrupt checkpoints, cancellation at each publication boundary,
and recovery near the original deadline. Negative cases must remain in the
request ledger with their attempted inputs and observed disposition. A valid
retry must not gain a fresh unlimited budget merely by restarting.
\citep{exec14_failure_design}

The current observations establish selected local behavior and expose a concrete
classification defect. They do not repair the implementation, rerun a live MCP
consumer, measure bank-scale recovery, or settle the previously observed
descriptor-replacement and FIFO limitations. Those claims require their own
tests and, for operating acceptance, a named consumer, target filesystem and
platform, identity authority, retention rules and declared service limits.
Public or generated fixtures can justify the prototype and demonstrate the
merit of later experiments before expensive bank trials are commissioned.
\citep{exec14_failure_design}
\label{exec14:failure-trials-end}


% SOURCE docs/release/2026-10-29/source/docs/writing-repair/2026-10-29/src/2388e126567f2b6e.tex CHAR 0
\section{What exact agreement across three interfaces establishes}
\label{exec16:semantic-parity}

The common operation service supplies a useful starting point for a banking
analyst, a command-line workflow and an assistant using MCP. Each should retain
the same question, assumptions, result and limitations. That objective needs a
precise boundary. Here the library entry is \texttt{Runtime.invoke}; a lower
level numerical helper is not automatically the same public operation. A CLI
exit code, an MCP protocol error and a canonical scientific disposition also
belong to different layers. Requiring identical bytes across those layers
would obscure their meanings rather than preserve them.

For a request $q$ that is admissible through all three negotiated surfaces,
let $L(q)$ denote the canonical library response, $C(q)$ the decoded CLI
response, and $M_s(q)$ and $M_t(q)$ the structured and decoded text responses
from MCP. The desired equality is
\[
 L(q)=C(q)=M_s(q)=M_t(q),\qquad q\in\mathcal Q_{\rm shared}.
\]
Equality includes context identity, operation, execution status, scientific
disposition, result kind, units, set topology, qualification, evidence and
error. A normalization that discards the error or reduces an unbounded set to
a convenient point would not meet this requirement. The shared domain is
itself negotiated: a generic local dispatcher can represent an unsupported
operation that the MCP server deliberately does not advertise.
\citep{exec16_parity_observed,exec16_parity_design}

\subsection{Actual local observations, with independent result meanings}

A deterministic local comparison exercised thirteen scientific requests
through the library, CLI and real Python MCP client/server over stdio. The
server negotiated protocol 2025-11-25 using SDK 1.27.0. All thirteen canonical
responses agreed exactly, including MCP's structured and serialized-text
copies. The cases comprised capabilities, graph syntax, a supplied-DAG
back-door expression, unsupported dependent sampling, a contradicted premise,
an unsupported estimation method, four deterministic sensitivity-set shapes,
weak-moment inversion, a changed-context failure and observation-clock
diagnostics. Three further lifecycle calls exercised report submission,
terminal status and cancellation after completion; together the selected
cases exercised all nine advertised operations. This is actual local execution
on generated inputs, not a mock transport or a statistical performance study.
\citep{exec16_parity_observed}

The equality check is useful because the expected meanings were also reviewed.
The four sensitivity examples solve $b^2\leq4$, $b^2+1\leq0$,
$-b^2-1\leq0$ and $-b^2+4\leq0$. Their sets are respectively
$[-2,2]$, the empty set, the whole line and
$(-\infty,-2]\cup[2,\infty)$. Their uncertainty is explicitly
nonprobabilistic. Matching copies of these objects cannot turn them into
sampling confidence regions. The moment example retains its conditional
Gaussian-moment interpretation. The graph example retains the supplied graph
and the limitation that syntax review does not prove causal truth. The clock
example flags both an observation available exactly at the decision time and
one without an availability time under the requested strict-before rule.
\citep{exec16_parity_observed}

For banking use, these distinctions matter when the same request is reused
in a batch report or presented by an assistant. The dependent-sampling and
unsupported-method examples returned a completed invocation with an
unsupported scientific disposition. The contradicted-premise example returned
a completed invocation with refusal. All three had CLI exit zero and
\texttt{isError=false} in MCP; none produced the requested causal estimate.
Changing the sign convention while retaining a moment artifact bound to the
old context instead produced a failed canonical response, CLI exit four and
\texttt{isError=true}. The preserved context hash and failure prevent a
response about one contrast from being relabeled as the opposite contrast.
This is agreement about the supported boundary, not a general guarantee that
every possible scientifically inconsistent context is recognized.
\citep{exec16_parity_observed}

The shared report job preserved computed, unsupported and refused input
declarations, while labeling its verification as schema-only. Its submission
response was a queued handle; the terminal state belonged inside the subsequent
job-status payload. Cancelling the completed job left its completed report
intact. A proposed consumer must display those states separately. A finished
report is not independent scientific approval of its contents, and a successful
cancel call cannot be presented as reversal of an already completed action.
This local case covers cancellation after completion; it does not replace
the distinct recovery, in-flight cancellation and target-platform evidence.
\citep{exec16_parity_observed}

\subsection{Capability negotiation and errors are part of the scientific contract}

The discovery request unimplemented in the tested service revision exposes the boundary of this equality statement. Later bounded discovery research does not retroactively change that recorded capability negotiation. Library and CLI returned the same canonical unsupported response,
with CLI exit zero. Calling the unadvertised discovery tool through MCP
produced JSON-RPC error $-32602$, ``Unknown causal tool.'' This follows the
adapter's separation between advertised operations and unimplemented canonical
requests. The protocol error says that this tool call is unavailable; it does
not say that a causal effect is nonidentified. A proposed client should retain
the negotiated capability set and resolve this availability question before
dispatch, using a common user-facing explanation while preserving the actual
wire error if a call is nevertheless made. It should not invent a scientific
response that the service never returned.
\citep{exec16_parity_observed,v12_mcp_legacy_tools,exec16_parity_design}

The invalid-schema case also behaved differently at the transport boundary.
An extra request property raised a library validation exception, produced
CLI \texttt{invalid\_request} with exit two, and produced an MCP tool error
without a canonical structured response. This is appropriate input rejection
for that case. The corrupt-replay case shows why a matching coarse error label
is not enough. A valid capability request was first saved under an idempotency
key; its local replay JSON was then deliberately malformed. The library raised
\texttt{JSONDecodeError}. The CLI returned \texttt{invalid\_request} with
exit two, while MCP returned ``Invalid local causal request'' and no canonical
response. The request had not changed. The failing bytes belonged to stored
execution state, so attributing the failure to the analyst's request is wrong.
This extends the earlier CLI failure finding with a real MCP observation.
The original replay bytes were restored after the isolated local trial.
\citep{exec16_parity_observed}

The source explains the attribution. Replay parsing occurs inside the runtime,
but JSON decoding errors inherit from \texttt{ValueError}; the CLI and MCP
adapters catch that broad exception family together with input-validation
errors. The proposed repair is to separate validated caller-input failures
from artifact, replay-state, worker and transport failures. For an already
validated request with a corrupt replay, return a failed/unavailable canonical
result with its context identity and a specific stored-state error. For an
unparseable incoming message, retain a transport/input error without fabricating
a scientific context. Preserve the corrupt record for operator investigation;
do not automatically discard an idempotency record and repeat a potentially
consequential action. This repair is specified, not implemented in the observed
runtime. \citep{exec16_parity_observed,exec16_parity_design}

\subsection{Validation ownership, identity and the remaining acceptance cases}

Validation responsibility cannot be assigned merely by noticing that a tool
advertises an output schema. In the inspected SDK, a handler returning an
explicit \texttt{CallToolResult} takes a path that returns that object before
the decorator's later structured-output validation branch. The local adapter
does explicitly validate its canonical response, and the runtime also checks
response semantics. A future refactor must preserve those checks and test
malformed responses on every return path. The successful local trials do not
establish that the SDK will perform validation which this specific code path
leaves to the adapter. \citep{exec16_parity_observed}

Execution identity needs an equally precise scope. The runtime binds selected
local sources, two schemas, Python and four numerical/validation dependencies.
Its identity does not itself include MCP or AnyIO versions or the separate
MCP tool-mapping contract. The current observations record SDK versions and
source identities separately. The proposed transport identity should bind
the negotiated protocol, tool map, SDK, concurrency dependency, launch profile
and any response normalizer as well as the scientific runtime. A common
scientific result with a changed transport contract then remains distinguishable
from an exact replay under the same complete interface.
\citep{exec16_parity_observed,exec16_parity_design}

The known next tests are concrete: valid requests with missing, malformed or
incompatible replay records; unknown tools and unsupported method identifiers;
schema and context mismatches; nonfinite or incorrectly typed responses;
divergent MCP text and structured content; changed SDK/tool-map identity;
and terminal versus in-flight job cancellation. Freeze the request and expected
layer of failure before each test, retain failed attempts, and compare the
whole result object. Extend the common decoder to each proposed consumer
language and then test against its actual caller. These engineering routes
are known, but their completion and external acceptance are not supplied by
the thirteen matched examples. \citep{exec16_parity_design}

The environment observation is also bounded. A six-second worker-thread probe
stalled under the managed sandbox; the same fixed local probe succeeded in
the approved execution environment outside it, where the MCP comparison ran.
This is an observed operating constraint, not a reason to treat sandbox and
production environments as interchangeable. HTTP, multiple operators,
authentication, remote retries, newer protocol revisions and bank-platform
performance remain separate qualification work. The immediate proposal can
justify testing these interfaces with generated and permitted public inputs;
actual consumer owners must later accept their deployed workflows. No transport
comparison can supply missing intervention truth, causal assumptions or useful
bank precision. \citep{exec16_parity_observed,exec16_parity_design}
\label{exec16:semantic-parity-end}


% SOURCE docs/release/2026-10-29/source/docs/writing-repair/2026-10-29/src/f5dcc88b9485dfca.tex CHAR 0
\section{Continuity across jobs, interfaces, and ownership}
\label{exec17:integration-continuity}

The local service already demonstrates several useful engineering ideas:
frozen requests, context-bound results, bounded workers, durable report
checkpoints, and matching responses across selected local interfaces. The next
proposal task is to state exactly what those observations warrant and what they
do not. A banking analyst may submit a portfolio-scale summary, inspect a
conditional treatment effect through a library, invoke the same request from a
command line, and review it in an agent-mediated workflow. Correctness requires
that a restart keeps the original scientific question, that each interface
preserves the meaning of results and refusals, that protocol and data schemas
are versioned independently, and that an institutional owner—not the author—
accepts a precisely identified release. These are related requirements; none
implies that another is already qualified.

The examples below concern a servicing report that summarizes canonical
responses for a frozen set of eligible accounts. “Actual” refers only to the
saved local executions on generated inputs. “Static” refers to source
inspection. “Authored reference” denotes a reproducible mathematical or finite
decision-table calculation. “Proposed” denotes a design awaiting its stated
acceptance tests. None is a bank result.~\citep{exec16_parity_observed,exec14_failure_observations}

\subsection{A durable report is not an arbitrary resumable estimator}
\label{exec17:continuity-jobs}

The report worker has a concrete recovery invariant. Let the submitted ordered
artifact identities be $(a_1,\ldots,a_m)$ and let $C_k$ be the ordered
identities already summarized. A valid checkpoint has
\[
 C_k=(a_1,\ldots,a_k),\qquad
 h_k=H(J,C_k,S_k),
\]
where $J$ binds the original request, execution identity, artifact hashes,
and declared budget, and $S_k$ is the corresponding ordered summary list.
Resume checks that the identities form this prefix, that there is one summary
per completed identity, and that the checkpoint digest matches. It refuses a
changed execution identity or checkpoint, caps starts at three, and computes
remaining wall time from the original creation time. Those checks protect a
declared report reconstruction. The worker summarizes response declarations
under a schema; it explicitly does not validate the causal claims inside them.
\citep{exec14_failure_implementation}

The saved fault trial killed a report worker after a checkpoint and later
observed a completed report with one entry per frozen input. A distinct
instrumented trial interrupted storage after content publication but before
artifact metadata; same-content retry completed the descriptor, while the
content-only orphan remained unlisted in the meantime. Another trial changed
checkpoint content and observed a resume refusal. These are actual local
observations with generated requests and controlled interruption hooks. The
hooks establish behavior at named write boundaries; they are not naturally
observed host or storage failures. The records cover one-host report
reconstruction, not arbitrary estimator continuation, multi-host fencing, or
exactly-once effects at an external bank service.
\citep{exec14_failure_observations,exec13_ops_faults}

There are at least three persistence transitions: a per-input checkpoint, the
content and descriptor writes for a content-addressed artifact, and the job's
terminal acknowledgement. A lock protects ordinary in-process transitions,
but a process crash releases it. If result bytes and metadata are durable and
the process fails before the terminal state is written, a later cancellation
may make the job terminally cancelled while retaining a report artifact that
was already published. This does not violate terminal immutability: the
terminal job state and an independently addressable immutable artifact are
different objects. The client must read the job handle and artifact reference
as separate authorities. Cancellation does not roll back durable bytes, sent
messages, or other effects. The current causalbank job operation applies only
to schema-summary reports; it is not MCP request cancellation or an arbitrary
scientific worker API. The legacy protocol itself distinguishes request and
task cancellation and requires parties to handle completion races.
\citep{v12_mcp_legacy_cancellation,exec14_failure_implementation}

An exact finite reference enumerates all sixteen ordered selections without
repetition of three checkpoint identities at lengths zero through three.
Exactly four are valid prefixes, including the empty prefix. It also separates
five lock-respecting cancellation placements from the crash window between
publication and terminal acknowledgement. In that window the proposed recovery
checks the exact input, report, and job-state identities, preserves the
immutable artifact, and lets an already-recorded terminal state win. If a
descriptor points to missing or mismatched bytes, the proposal refuses
automatic repair and retains the evidence for operator reconciliation. These
counts are combinatorial checks of the stated finite model, not failure rates
or measurements of the implementation.
\citep{exec17_continuity_reference}

The deadline creates a separate static concern. Current resume logic computes
\[
 r_{\rm wall}=t_{\rm created}+B-t_{\rm now},
\]
where $B$ is the original budget and both $t$ values are wall-clock readings.
With $t_{\rm created}=1000$, $B=60$, and fifty seconds of real elapsed time,
the formula returns ten seconds. If eighty seconds have actually elapsed while
the clock has moved backward ninety seconds, the same formula returns seventy
seconds, even though the real elapsed budget has expired. This is a
mathematical counterexample to the formula under a wall-clock adjustment, not
an observed clock failure. A monotonic clock is the proper within-process
measure; a reboot-safe cross-process deadline needs an explicit trusted time
source or a conservative failure policy, because a monotonic epoch is not
portable across restarts. The accepted design must choose and test that
policy rather than silently promise both properties.
\citep{exec14_failure_implementation,exec17_continuity_reference}

The remaining local acceptance matrix should freeze tests for every worker
boundary: before and after each checkpoint; content before metadata; metadata
before terminal acknowledgement; cancellation before work, during work,
after publication, and after completion; repeated resume requests; stale
leases; process death; malformed or truncated state; changed code, dependency,
schema, artifact, or launch identity; deadline expiry; backward and forward
clock changes; and duplicate external-effect attempts. Each trial must compare
the full request identity, ordered prefix, artifact bytes and descriptor,
terminal state, attempt count, and emitted effects. Expected refusals belong in
the denominator. A fault harness may force the process boundary, but the
target platform owner must later run the frozen protocol on its declared
storage, scheduler, hardware, clock, and recovery topology. Current external
input contracts leave the platform operator and recovery lead unassigned.
\citep{exec17_continuity_contract}

\subsection{One semantic operation domain across three surfaces}
\label{exec17:shared-operation}

Let $\mathcal Q_L$, $\mathcal Q_C$, and $\mathcal Q_M$ denote requests admitted
by the library, CLI, and MCP tool map. The defensible parity target is the
negotiated intersection
\[
 \mathcal Q_* = \mathcal Q_L\cap\mathcal Q_C\cap\mathcal Q_M,
 \qquad
 L(q)=C(q)=M_{\rm structured}(q)=M_{\rm text}(q),\quad q\in\mathcal Q_*.
\]
Equality means the same canonical response after transport decoding: request
and context hashes, operation, invocation status, scientific disposition,
units, estimate or identified set, topology, uncertainty meaning, evidence,
qualification, and typed error. It does not mean identical JSON-RPC framing,
CLI exit code, or tool-level error wrapper. Those belong to different layers.
The MCP server must expose only operations in its actual mapping; the generic
runtime may return a canonical unsupported result for an operation the tool
map does not advertise.~\citep{exec16_parity_design,exec16_parity_observed}

The actual saved parity suite ran thirteen generated scientific requests and
three lifecycle comparisons through the local library, CLI, and MCP stdio
client/server. Complete canonical results agreed in the tested common domain;
all nine advertised operations were exercised. Separate adverse tests observed
an unadvertised tool, schema rejection, and a valid request whose corrupted
stored replay was mislabeled as invalid caller input by broad adapter error
handling. The latter is a concrete error-origin defect, not a causal refusal.
The suite establishes these selected cases for the recorded Python and SDK
versions. It does not establish full parity over every schema branch, a newer
SDK, another language, HTTP transport, or a bank consumer.
\citep{exec16_parity_observed}

A shared-operation design should have one versioned registry that binds the
canonical request schema, response validator, method identifier, common
capability, and typed failure origin. The library runtime owns dispatch and
scientific context checks. The CLI owns bounded input parsing, command-name
agreement, JSON serialization, and documented exit mapping. MCP owns the
negotiated tool list, wire protocol, and error wrapper. Each adapter must call
the same canonical runtime and validate the same response object before
encoding it. A protocol-level “unknown tool” remains a capability error;
canonical unsupported analysis remains a scientific disposition; malformed
input remains an input error; corrupt replay remains stored-state failure;
and worker death remains worker failure. A successful transport is not a
scientific success. This design preserves semantic ownership while preventing
each adapter from growing an independent operation implementation.
\citep{exec16_parity_design,exec14_failure_design}

For an account-value request, changing the treatment sign, the thirty-day
horizon, USD units, or independent-account sampling law changes the scientific
target even if the JSON schema version remains unchanged. The context identity
must therefore change, and a previous moment or audit artifact bound to the
old context must be rejected. Conversely, changing a CLI wrapper without
changing the target should require adapter conformance evidence without
misrepresenting it as a new causal estimand. A reproducible authored example
uses a stipulated effect set $[-2,6]$ USD per account and a cost of USD $3$.
For $1000$ accounts the net-value set is
\[
 1000\bigl([-2,6]-3\bigr)=[-5000,3000]\text{ USD}.
\]
Converting both effect and cost to thousand-dollar units and then converting
the result back leaves that set unchanged. Converting only one quantity does
not. This exact interval arithmetic demonstrates why response equality must
preserve units and sign; it is not an estimate or a profit claim.
\citep{exec17_continuity_reference}

The repair sequence is concrete: define the registry and common domain;
separate caller, protocol, stored-state, artifact, worker, and transport
failures; freeze request/error pairs before running; test full-object equality
on every adapter branch; include changed context and artifact mismatches;
then add actual maintained consumers one at a time. A consumer's decoder must
preserve null endpoints, disconnected sets, large integers, explicit refusal,
and uncertainty labels. Current inputs identify consumer repositories and
platform/security owners but provide neither their actual decoder contracts
nor signed workflow acceptance. Language adapters can be specified now;
their external acceptance cannot be authored locally.
\citep{exec16_parity_design,exec17_continuity_contract}

\subsection{Version the scientific contract separately from its transport}
\label{exec17:versions}

“Version” refers to several identities. The proposed compatibility manifest
records these independently:
\begin{center}
\begin{tabular}{P{0.19\linewidth}P{0.28\linewidth}P{0.43\linewidth}}
\hline
Identity & Example & Change that requires review\\
\hline
Scientific context & Target, sign, horizon, units, assumptions, clocks, sampling & Changes estimand or interpretation; invalidate context-bound evidence\\
Request/response schema & \texttt{request/1.0.0} & Changes fields, constraints, or representation; migrate and run schema tests\\
Operation contract & \texttt{local-runtime/1} & Changes method domain or result/refusal semantics; run operation conformance\\
Artifact schema & \texttt{causalbank.report/1} and input-role schemas & Changes persisted object meaning; migrate or refuse stale artifacts\\
Wire protocol & MCP \texttt{2025-11-25} or \texttt{2026-07-28} & Changes framing, negotiation, tool result, or session behavior; qualify client/server pair\\
Optional extension & Namespaced negotiated task or other extension & Changes opt-in behavior; require explicit capability agreement\\
\hline
\end{tabular}
\end{center}

The current source pins local request/response schemas at 1.0.0 and reports
operation version \texttt{local-runtime/1}; reports use their own artifact
schema. The MCP adapter is pinned to the 2025-11-25 initialize/initialized
protocol and Python SDK 1.27.0. The pinned 2026-07-28 specification instead
puts version metadata on each request and defines server discovery. Modern
servers must implement discovery; a modern client may call another request
directly and handle a version error. This differs from the recommended probe
for a client that must distinguish modern and legacy servers. The
installed SDK snapshot declares the legacy protocol as its latest version;
modern support is not implemented or trialed here. These are audited source
facts and bounded runtime observations, not a claim that the newer standard is
unsupported by all SDKs.~\citep{v12_mcp_legacy_lifecycle,v12_mcp_versioning,v12_mcp_basic,v12_mcp_stdio,exec16_parity_observed}

The protocol transition rule is subtle. A dual-era stdio client probes
\texttt{server/discover}. A discovery result means modern; the client chooses
a mutually supported version. A recognized modern protocol-version error
also means modern: retry with one of its advertised supported versions or
report incompatibility, but do not fall back to legacy initialization. A
different error or timeout identifies the legacy era, so a dual-era client
tries the legacy handshake. The fallback cannot depend on one specific
unknown-method error code because legacy implementations may return several
codes or no response. This is a protocol decision, not a scientific
operation.~\citep{v12_mcp_stdio,v12_mcp_versioning}

The authored decision table covers a discovery result with and without a
common version, a recognized unsupported-version error, another recognized
modern error, two ordinary errors, timeout, and a modern-only client. It
retains the no-downgrade rule for every recognized modern error. The table is
not a wire parser and no simulated row establishes actual SDK behavior. Before
advertising dual-era support, a real client/server matrix must test malformed
discovery bodies, empty and contradictory version lists, advertised-version
retry, recognized errors, unknown errors, timeout, cache lifetime, and server
restart. It must also compare the exact canonical response and tool capability
map across each supported protocol.~\citep{exec17_continuity_reference,exec17_continuity_contract}

Compatibility must not silently widen the scientific request. A protocol
version negotiates message behavior; it cannot decide whether an interval is
a confidence region, convert a posterior set to a sensitivity region, or
repair an ambiguous treatment. A schema version does not prove assignment
truth, temporal ordering, or graph-class validity. The adapter should reject
unknown mandatory semantic fields, support only explicit migrations, and
retain original request bytes and version identities for audit. Unrecognized
scientific semantics require refusal or a new reviewed operation contract,
not best-effort coercion.
\citep{v12_mcp_basic,exec16_parity_design}


\label{exec17:continuity-end}


% SOURCE docs/release/2026-10-29/source/docs/writing-repair/2026-10-29/src/1b396b1d3bc753ab.tex CHAR 0
\hypertarget{a-java-discovery-boundary-that-retains-the-question}{%
\section{A Java discovery boundary that retains the
question}\label{a-java-discovery-boundary-that-retains-the-question}}

\label{exec13:java-adapter}

The Java comparison motivates a concrete first integration: a pinned
Tetrad process for a restricted PC discovery request, with a separate
admission path for FGES. The initial request concerns complete
continuous observations under an explicitly assumed IID linear Gaussian
model with causal sufficiency, acyclicity, faithfulness and no unmodeled
selection. These restrictions identify a reviewable prototype. They do
not imply that a table of banking observations satisfies them. Clustered
customers, repeated account histories, missing-data correction,
latent-variable discovery and background-knowledge constraints require
other profiles. The adapter is a proposal contract; no Java engine or
new transport endpoint has been implemented here.
\citep{exec13_java_adapter_contract,exec13_java_tetrad}

The request retains the downstream target even though discovery does not
compute it. A funding analyst's population, treatment contrast, outcome,
time horizon, units and sampling declarations remain attached to the
graph. They prevent a later identification step from silently changing
the question. The generic operation contract already provides these
fields for effect operations. The Java discovery profile deliberately
requires that context as well. A standalone graph request without a
downstream treatment and outcome needs its own reviewed context
extension; inserting arbitrary variable names would create a fictitious
causal question. \citep{exec13_java_adapter_contract}

The logical interface is the same across the proposed library
invocation, \texttt{causalbank operation graph.discover} command and
\texttt{causal\_graph\_discover} MCP tool. Each transports a versioned
canonical request and receives the same typed result. A complete
scientific context digest accompanies a broader request digest
containing input hashes, engine revision, algorithm settings, adapter
version and resource limits. An interface that changes USD into
thousands of USD, reverses the treatment contrast, changes the sampling
law or orients an ambiguous edge has changed the scientific request,
even if its JSON remains valid. The reference serialization examples
demonstrate those identity changes on local objects; they are not
executions of a server or MCP client.
\citep{exec13_java_adapter_contract,exec13_java_adapter_oracles}

\hypertarget{the-graph-object-preserves-a-class-of-possible-interventions}{%
\subsection{The graph object preserves a class of possible
interventions}\label{the-graph-object-preserves-a-class-of-possible-interventions}}

Use opaque variable identifiers and a bijection into engine symbols.
Store the original column order and its permutation separately from the
engine's order. Every variable, including an isolated node, travels
through the conversion. Human-readable labels, availability times and
units remain metadata with their own scientific meaning; names do not
supply an orientation rule. For every pair, retain its two endpoint
marks. When the serialization orders the pair lexically, it swaps the
marks with the nodes. Thus \(R\to B\), expressed as a tail at \(R\) and
an arrowhead at \(B\), agrees with the reversed pair representation
having an arrowhead at \(B\) and a tail at \(R\). Reversing only the
node order changes the graph.
\citep{exec13_java_tetrad,exec13_java_adapter_oracles}

A Gaussian example illustrates why an unoriented link must survive this
boundary. Consider three standardized variables \(X,Y,Z\) with
covariance \[
 \Sigma=\begin{pmatrix}
 1&4/5&2/5\\
 4/5&1&1/2\\
 2/5&1/2&1
 \end{pmatrix}.
\] For a banking interpretation, \(X\) can denote a rate concession,
\(Y\) a balance response and \(Z\) a funding outcome, all in stipulated
standardized units. The covariance obeys \(X\perp Z\mid Y\). The
undirected path \(X-Y-Z\) represents three acyclic orientations without
an unshielded collider. The orientation \(X\to Y\leftarrow Z\) is
excluded because it would require a different independence pattern under
faithfulness. \citep{exec13_java_adapter_oracles}

For each retained orientation, regress a variable on its proposed
parents using the same \(\Sigma\), and assign the residual variance as
its independent Gaussian noise variance. In matrix notation, if
\(W_{ij}\) is the coefficient from source \(i\) to target \(j\), then \[
 \Sigma=(I-W^{\mathsf T})^{-1}\Omega(I-W)^{-1},
\] where \(\Omega\) is the diagonal noise covariance. The exact
reference calculation verifies this equality for all three models. In
the forward chain, the coefficients are \(4/5\) and \(1/2\); intervening
on \(X\) changes \(Z\) by \(2/5\) per unit. In the fork and the reverse
chain, \(X\) has no directed path to \(Z\), so the effect is zero. The
same observational covariance therefore supports the effect multiset
\(\{2/5,0,0\}\). \citep{exec13_java_adapter_oracles}

Returning the forward chain alone would select the positive effect by an
unrecorded convention. Returning the interval \([0,2/5]\) would
introduce values such as \(1/5\) that none of these three specified
linear Gaussian models produces. That interval is a hull, not the
discrete effect set, and neither object is a confidence set without a
separate statistical construction. The discovery response consequently
retains the CPDAG and returns no effect estimate, effect region or
uncertainty level. A later identification operation must state how it
handles the whole class and the additional functional assumptions.
\citep{exec13_java_adapter_contract,exec13_java_adapter_oracles}

Checking endpoint syntax is insufficient to award the CPDAG label. A
graph can contain only tails and arrows and still fail to represent a
completed Markov equivalence class. The existing reference validator
does not recognize all such cases. The proposed Java boundary therefore
requires an independently reviewed graph-class check before publishing a
completed CPDAG. A circle, bidirected conflict or invalid completion in
the initial route is an invalid-engine-output failure, with the raw
graph retained for diagnosis. It is not silently coerced into a DAG.
Background knowledge is initially excluded because orienting an
observational CPDAG using additional restrictions requires an
appropriate constrained-graph representation and provenance.
\citep{exec13_java_adapter_contract}

\hypertarget{admission-must-cover-the-failure-path-as-well-as-the-constructor}{%
\subsection{Admission must cover the failure path as well as the
constructor}\label{admission-must-cover-the-failure-path-as-well-as-the-constructor}}

The proposed PC profile sets the conditioning depth to the inspected
unbounded value, enables the stable skeleton option, selects the
separating-set collider rule, disallows bidirected conflicts, and
prevents directed cycles. It supplies empty background knowledge and a
caller-declared test level. These choices make the attempted operation
explicit. A user-requested depth cap is initially unsupported;
exhausting a runtime budget returns an incomplete job, not a graph
advertised as the result of the completed search.
\citep{exec13_java_adapter_contract,exec13_java_adapter_source}

The numerical test profile is equally specific: complete data, no
shrinkage, zero ridge and no pseudoinverse fallback, with all settings
frozen before search. The inspected Fisher-Z source exposes these
alternatives separately. Automatically activating regularization after a
singularity would change the test while keeping the same apparent method
name. That may be useful in a different profile, but it needs a separate
calibration argument. The initial profile requires at least
\(n\geq p+2\) complete observations for \(p\) variables so that even the
largest possible conditioning set has positive Fisher-Z degrees of
freedom. This is only an arithmetic prerequisite; it does not guarantee
useful power or stable covariance inversion.
\citep{exec13_java_adapter_contract,exec13_java_adapter_source}

As the Java comparison showed, the inspected \texttt{checkIndependence}
path converts a singular-matrix exception into a dependence result with
\(p=0\). The proposed wrapper instead uses the exception-preserving
\texttt{getResult} path and derives a decision only from a valid result.
Every invocation receives a query identity and a terminal record. Before
propagating an exception, the wrapper sets a job-level failure flag.
Publication then checks that every issued call has a valid terminal
record and that the flag is clear. This matters because search code may
catch interruptions internally: receiving a graph object alone does not
prove that all its requested tests completed.
\citep{exec13_java_tetrad,exec13_java_adapter_source,exec13_java_adapter_contract}

That wrapper remains an implementation prerequisite. Its full
independence-test interface, ancillary methods and interaction with the
search implementation must be inspected and compiled; calling one method
in a design document is not sufficient. Numerical rank checks before
dispatch help detect duplicate accounting columns, but per-query records
remain necessary for failures that arise during the search. A zero test
probability cannot be reinterpreted as maximal causal confidence. If the
required failure provenance cannot be obtained, the route is unavailable
in this runtime. \citep{exec13_java_adapter_contract}

Parallelism and deadlines also belong to admission. The stable FAS path
uses a parallel stream and requires thread-safe records. FGES's separate
worker setting does not control that path. The inspected PC search
delegates the skeleton phase to FAS without demonstrating an end-to-end
deadline guarantee from PC's timeout setting alone. The proposed parent
process therefore enforces the whole-job deadline, memory and output
budget. An engine's cooperative timeout is supplementary. A cancellation
or a process death preserves the attempt's failure record and withholds
its graph.
\citep{exec13_java_adapter_source,exec13_java_adapter_contract}

\hypertarget{units-unsupported-requests-and-the-next-executable-comparison}{%
\subsection{Units, unsupported requests and the next executable
comparison}\label{units-unsupported-requests-and-the-next-executable-comparison}}

The signed funding model in the Java survey provides a useful
independent conversion reference. If \(x'_i=d_i x_i\), the source-row
coefficients must transform as \(W'_{ij}=d_jW_{ij}/d_i\). Rescaling the
rate by 100, balances by 1000 and cost by 10, then applying the
correspondingly rescaled rate intervention, must recover the original
responses after converting units back. The exact reference does so. It
also verifies that an unchanged absolute coefficient threshold can
select a different graph after rescaling. This does not give
coefficients to the PC adapter, whose result is graph-only; it supplies
an admission example for any later weighted route.
\citep{exec13_java_adapter_oracles,exec13_java_examples}

A valid request outside the initial scientific scope receives a
completed unsupported response with its context preserved and a specific
alternative. Repeated account observations point toward a separately
reviewed temporal or cluster profile; latent-variable assumptions point
toward an appropriate PAG route; missing values require an explicit
observation mechanism and correction profile. An uninstalled or
unqualified runtime receives an unavailable response. Neither condition
is called nonidentification, which would be a stronger scientific
conclusion about an effect under a model. Malformed inputs and failed
computations carry structured errors, with no effect or graph result
masquerading as a completed calculation.
\citep{exec13_java_adapter_contract}

The first prototype should use a fixed local Java process so its
dependency and runtime identity can be measured. The process boundary
does not remove the inspected GPL conditions. The actual compiler
target, runtime vendor and version, dependency artifacts, operating
system and thread configuration must be locked for the chosen route; the
source's compiler setting and bridge README cannot substitute for that
environment check. Result publication is atomic after validation. A
crashed search is restarted as a new retained attempt under the same
immutable request, with no promise of resuming an internal search state.
If completion wins a cancellation race, the immutable published result
remains the observed outcome.
\citep{exec13_java_tetrad,exec13_java_adapter_contract}

The executable comparison should begin with independent oracles: the
three-model Gaussian class, a collider and an isolated node, the
depth-three separator, signed matrices and unit changes, singular input,
and injected failures. Only after these pass should a prospectively
specified generated-sample study estimate graph error and resource use
under the declared assumptions and their violations. It should retain
every attempted request, including timeouts and refusals, and compare a
matched baseline under the same budgets. These steps test whether the
proposed boundary deserves a prototype claim. Expensive bank experiments
become relevant when that bounded evidence exposes a decision-relevant
uncertainty that public or generated data cannot settle.
\label{exec13:java-adapter-end}

\hypertarget{encoding-algorithm-settings-without-changing-the-shared-schema}{%
\subsection{Encoding algorithm settings without changing the shared
schema}\label{encoding-algorithm-settings-without-changing-the-shared-schema}}

The existing discovery request schema permits only an algorithm
identifier and algorithm version in its parameter object. The proposed
Java profile therefore places the detailed settings in a separate
immutable configuration artifact. The request supplies its opaque
handle, hash, media type, schema version and role, alongside the data
artifact. Resolving that artifact recovers the test level, conditioning
depth, missingness policy, numerical controls and background-knowledge
choice. The full request identity includes its hash, so changing a
setting creates a different request even when the downstream scientific
context remains identical. This preserves compatibility with the shared
schema while retaining the scientifically relevant settings.
\citep{exec13_java_adapter_contract}

The adapter's admission check must require exactly one configuration
artifact of the reviewed version, verify its bytes and recognized
fields, and check that the declared algorithm version admits that
profile. Generic schema validation alone cannot establish those
conditions. The worked wire records separately represent an unavailable
runtime, an unsupported sampling regime, a failed invocation and a graph
response from an independent population reference. Their data fixture is
deterministic, and their graph is not a fitted result. These records
make the proposed response contract concrete without treating a valid
message as evidence of an executed engine. \label{exec13:java-wire-end}


% SOURCE docs/release/2026-10-29/source/docs/writing-repair/2026-10-29/src/d1cfda61285f1f9e.tex CHAR 0
\section{A native forest boundary that preserves the banking question}
\label{exec13:grf-adapter}

A servicing analyst asks whether making a specified offer changes
twelve-month net receipts for eligible customers. The proposed treatment
is the offer, coded one, against no offer, coded zero. The outcome
includes the offer's cost and is measured in dollars per customer over
the stated horizon. A heterogeneous-effect forest can help estimate this
contrast under a declared identification model. Its C++ implementation
cannot decide whether the recorded offer was randomized, whether all
confounders were measured, or whether the analyst actually wants the
mean over customers, branches, or treated customers. Those choices must
precede engine dispatch and survive it.

The proposed first grf integration uses its maintained R interface in a
pinned local process, which in turn invokes the native C++ core. Direct
native embedding is a separate alternative. Inspection covers commit
\texttt{5bee99b51471}, dated 30 April 2026, whose package description
declares version 2.6.1. The snapshot therefore precedes the literature
cutoff; its October acquisition date is separately retained. The design
is not an installed adapter. The executed examples below are
project-owned arithmetic references; they neither train grf nor measure
its statistical performance
\citep{exec13_grf_runtime,exec13_grf_design,exec13_grf_oracle}.

\subsection{From an offer question to a local moment}

Let \(X\) contain information available before the offer,
\(W\in\{0,1\}\) denote assignment, and \(Y(w)\) denote the potential
twelve-month net receipt. With consistency, no interference, conditional
exchangeability and positivity, the target is \[
 \tau(x)=E[Y(1)-Y(0)\mid X=x],\qquad
 \theta=E[\tau(X)].
\] The outer expectation defines the customer population. An observed
treated subpopulation or an overlap-weighted population gives a
different target. The initial proposed profile accepts complete
observations from independent customers, a binary treatment, one fixed
outcome and a fixed horizon. It excludes repeated decision histories,
unmeasured confounding, interference, survey weighting and latent-graph
ambiguity that has not been resolved for this estimand. These are
restrictions of the proposed profile, not assertions that the package
has no other functionality.

Write \(m(x)=E[Y\mid X=x]\) and \(e(x)=E[W\mid X=x]\). The inspected R
wrapper either obtains nuisance predictions or uses supplied
predictions, then passes \(Y-\widehat m(X)\) and \(W-\widehat e(X)\) to
the native training boundary. A direct call that passes the original
vectors does not replicate that boundary. The causal binding selects the
instrumental trainer and sets the instrument coordinate to the treatment
coordinate. Consequently the inspected native prediction is a ratio of
locally weighted covariances, with the treatment variance in its
denominator \citep{exec13_grf_rpath,exec13_grf_native}.

For normalized nonnegative forest weights \(a_i(x)\), let
\(\widetilde Y_i\) and \(\widetilde W_i\) denote the residualized
values. The local calculation has the form \[
 \widehat\tau(x)=
 \frac{\sum_i a_i(x)(\widetilde Y_i-\overline Y_a)
                         (\widetilde W_i-\overline W_a)}
      {\sum_i a_i(x)(\widetilde W_i-\overline W_a)^2},
 \quad
 \overline Y_a=\sum_i a_i(x)\widetilde Y_i,
 \quad \overline W_a=\sum_i a_i(x)\widetilde W_i.
\] In the exact reference, equal weights on residual outcomes
\((-2,-1,1,2)\) and residual treatments \((-1/2,-1/2,1/2,1/2)\) give
slope three. Constant residual treatment gives zero denominator and the
reference refuses the calculation. Its numerical threshold is an
explicit arithmetic rule, not evidence that positivity holds in a bank
population. The native source divides by its local first-stage quantity;
the adapter must inspect finite outputs and retain weak-support
diagnostics instead of silently replacing a failed division with zero.

Honest tree splitting and out-of-bag prediction address particular reuse
of observations. They do not establish the causal premises or turn
customer rows from a common branch shock into independent samples.
Generalized forest theory also imposes regularity, subsampling, smooth
local identification and variance conditions; a named package and a
successful fit do not verify them \citep{grf2018}. The proposed first
profile retains the R nuisance path and its out-of-bag predictions.
Supplying an independently estimated nuisance vector requires a later
profile with explicit training identities, alignment and sampling-unit
split rules. A known randomized propensity may be supplied only with the
corresponding assignment declaration; it is never inferred from a
convenient constant.

\subsection{Why the average effect needs its own calculation}

The average of fitted conditional effects is not interchangeable with
the package's default augmented average-effect estimator. In the
inspected binary path, write \(\widehat\tau_i\) for the out-of-bag
conditional prediction. Its score is \[
 \Gamma_i=\widehat\tau_i+
 \frac{W_i-\widehat e_i}{\widehat e_i(1-\widehat e_i)}
 \{Y_i-\widehat m_i-\widehat\tau_i(W_i-\widehat e_i)\}.
\] This equals the familiar two-outcome-regression expression after
setting \(\widehat\mu_{0i}=\widehat m_i-\widehat e_i\widehat\tau_i\) and
\(\widehat\mu_{1i}=\widehat m_i+(1-\widehat e_i)\widehat\tau_i\). The
equality is algebraic. Consistency and uncertainty still require
appropriate nuisance estimation, overlap and sampling assumptions
\citep{exec13_grf_rpath}.

The reference uses four tuples
\((Y,W,\widehat m,\widehat e,\widehat\tau)\): \((0,0,1,1/4,2)\),
\((1,0,1,1/2,2)\), \((3,1,2,1/2,2)\), and \((5,1,2,3/4,2)\). Their
scores are \(8/3,0,2,16/3\). The mean score is \(5/2\), whereas the mean
conditional prediction is two. The C++ residual-score calculation and an
independent rational Python calculation using the two-regression formula
agree. This is an exact conversion check with deliberately supplied
numbers, not fitted evidence for double robustness or coverage
\citep{exec13_grf_oracle}.

For target weights \(q_i\), the corresponding estimator is
\(\widehat\theta=\sum_iq_i\Gamma_i/\sum_iq_i\). Choosing \(q_i=1\)
targets a customer-weighted population; choosing inverse branch sizes
gives each observed branch equal total weight. A branch with two
customer scores of ten dollars and a branch with four scores of forty
dollars produce a customer mean of thirty dollars and an equal-branch
mean of twenty-five dollars. A planning cohort of ten thousand customers
turns the five-dollar difference into fifty thousand dollars in a naive
projection. That calculation exposes the target change; it does not
justify applying an equal-branch mean to an arbitrary customer
portfolio.

The R interface documents that equal cluster weighting changes both
training sampling and average-effect weights, and disallows user sample
weights in that mode. Replacing the final mean alone cannot reproduce
training under that option. A native port would have to preserve both
mechanisms. The initial independent-customer profile therefore rejects
cluster and weight options; the independent-cluster extension is
specified for later qualification and must name whether customers or
clusters define its target \citep{exec13_grf_rpath}.

\subsection{Sampling uncertainty is not a numerical decoration}

For independent clusters indexed by \(g\), define
\(S_g=\sum_{i\in g}q_i(\Gamma_i-\widehat\theta)\), \(Q=\sum_iq_i\), and
let \(G\) count positive-weight clusters. The inspected all-sample
augmented average-effect path computes the variance expression \[
 \widehat V_{\rm cl}=\frac{G}{G-1}\frac{\sum_g S_g^2}{Q^2}.
\] This specifies a computation. Asymptotic validity needs an
appropriate independent-cluster experiment, sufficient effective
clusters and nuisance conditions; no fixed small cluster count
guarantees it. Dependence shared across branches through a calendar
factor lies outside that law.

For the six scores above, treating all rows as independent gives
variance forty dollars squared. Preserving two branch clusters while
retaining the customer target gives \(1600/9\) dollars squared. Giving
the two branches equal total weight gives variance 225 dollars squared
around the changed target of twenty-five dollars. The standard errors
are approximately 6.325, 13.333 and 15 dollars. Two clusters make a
useful arithmetic counterexample and supply no empirical
confidence-coverage qualification. A single positive-weight cluster is
refused. Zero total weight is separately refused
\citep{exec13_grf_oracle,exec13_grf_rpath}.

Conditional prediction variance is another quantity. The native
prediction strategy estimates it using grouped trees and a
variance-debiasing calculation; the binding copies the variance values
into \texttt{variance.estimates}. The R average-effect operation instead
returns \texttt{std.err}, already square-rooted. Squaring the wrong
field or taking its square root twice changes the units. The proposed
adapter keeps the raw field identity, method and units, then creates any
reported interval only under a declared inference profile. Pointwise
conditional intervals cannot become a simultaneous customer band or an
interval for their average through serialization. Likewise
\texttt{excess.error}, a finite-forest error diagnostic, is not a
substitute for sampling uncertainty
\citep{exec13_grf_native,exec13_grf_rpath}.

\subsection{Memory ownership, coordinates and model persistence}

The inspected \texttt{Data} class is a read-only view of a borrowed
contiguous column-major buffer. Element \((r,c)\) is read at offset
\(cn+r\), where \(n\) is the number of rows. Constructing the view from
a temporary vector does not transfer ownership. The host must keep the
allocation alive until all training or prediction threads have finished;
it must validate shape, index ranges and multiplication overflow before
the call. A three-row, two-column example has columns \((1,2,3)\) and
\((10,20,30)\). Correct layout reads twenty at zero-based coordinate
\((1,1)\); passing row-major bytes to the same index formula reads
three. Neither value is a memory error, so a mere successful call cannot
detect the scientific corruption
\citep{exec13_grf_memory,exec13_grf_oracle}.

Role coordinates must remain separate from predictor coordinates. The
native setters mark outcome, treatment, instrument and weight columns as
disallowed split variables. The host records the ordered covariate
dictionary, row identifiers, role indices and training/query
distinction. It cannot infer those roles from whichever columns happen
to be numeric. Missing values in covariates are supported by inspected
package paths, but are excluded by the initial profile so that an
undocumented imputation or missingness assumption is not introduced.
Conversion to double precision must preserve identifier strings
separately and refuse nonfinite numerical values rather than silently
dropping affected customers.

Model serialization also has ownership semantics. The inspected Rcpp
helper iterates by moving the unique pointer to each native tree into
its serialized representation. The header explicitly says the original
forest cannot be used after that operation. A direct embedding that
serializes and then predicts on the consumed forest is therefore
invalid. Either serialize an explicitly supported copy, or deserialize
the saved representation into a new object. Retain the engine commit,
ordered feature map, original training-data identity, nuisance vectors,
options and format version beside the forest. The selected source does
not establish a stable C++ ABI or a portable standalone model format.
These are admission questions for a native route, not problems solved by
choosing a file extension \citep{exec13_grf_memory}.

\subsection{A concrete request and result contract}

The operation is \texttt{estimate.effect}. Its generic parameter schema
accepts an engine identifier and estimator identifier, so the proposed
profile uses \texttt{grf-r-native-2.6.1} and
\texttt{binary-iid-aipw-all-v1} there. Detailed settings reside in an
immutable configuration artifact, carried through the existing artifact
list. That artifact binds the exact source commit, treatment contrast,
outcome/horizon, row and column roles, seed, thread count, tree options
and target. The adapter validates its own profile before dispatch;
generic JSON acceptance alone does not validate these scientific
restrictions \citep{exec13_grf_design}.

Library, CLI and MCP surfaces must resolve identical authorized data and
configuration identities. The context digest retains the estimand,
population, units, information clock, assumptions, sampling declaration
and graph. The full request identity additionally binds data bytes,
configuration, engine environment and numerical settings. Changing the
target population or turning branch weighting on produces a different
request; transport names cannot cause reuse of a previous result. Graph
endpoint marks and unresolved sets remain intact. A forest is not a
graph-identification service: if a retained model class gives different
effects and the query has no accepted identification argument, this
profile cannot silently select a graph.

Before runtime admission, discovery of this engine returns unavailable.
A recognized but unsupported sampling regime, continuous treatment,
dynamic policy request or unresolved graph requirement returns
unsupported with an actionable explanation. Nonfinite output, failed
subprocess execution or an invalid model artifact is a failure; no
finite effect accompanies it. A reviewed identification failure can
return not-identified with its evidence, but the adapter's inability to
analyze a model is not an impossibility proof. An overlap alarm produces
a refusal under the declared profile. Trimming or switching to an
overlap population requires a newly reviewed target and request.

After engineering admission, a computed scalar estimate keeps its dollar
unit, signed offer-minus-no-offer contrast and target identifier. A
diagnostic conditional-prediction artifact retains query-row coordinates
and explicitly denotes conditional means, not individual counterfactual
differences. The initial engineering result carries execution-only
qualification and no confidence claim. Conditional inference needs its
own reviewed uncertainty construction and evidence. A mathematically
unbounded or disconnected result from another method must be returned
with its topology; this forest profile does not approximate it by a
large finite interval.

\subsection{Deployment choices and the next informative comparison}

The package description requires R at least 3.5.0, C++17 and GNU make,
links Rcpp and RcppEigen, and imports additional packages including
Matrix, DiceKriging, lmtest and sandwich. These declarations are
minimum/source requirements, not a tested environment lock. The proposed
R process must record the installed R and dependency versions, source
and binary hashes, compiler and flags, thread settings and runtime
options. The inspected native CMake file collects both source and test
files into an executable; its presence does not supply a ready
production shared library or a reviewed installed ABI. Package metadata
declares GPL-3 and the inspected source headers state GPL version 3 or
later. Preserve file and dependency terms and review distribution
obligations before shipping; a subprocess boundary alone settles neither
licensing compatibility nor distribution duties
\citep{exec13_grf_runtime}.

The parent process bounds wall time, memory and output size and
publishes only a complete validated result atomically. Standard output
is the protocol channel; diagnostics have a bounded separate log. A
crash or timeout leaves a typed failure and its input identities. No
mid-tree recovery is promised. Replaying an immutable request starts
from its recorded seed and environment; bitwise equality across
different compilers or thread settings remains a measured property. One
worker per isolated R process is the conservative initial choice because
the inspected binding also writes global runtime context. Concurrency
inside a shared embedding requires its own review.

The proposal recommends the R/native route first because it retains the
inspected nuisance, score and aggregation paths. Direct native embedding
may later reduce interface overhead or fit a host's deployment
constraints, but would transfer those statistical and ownership
responsibilities to the host. A full algorithm port has the greatest
maintenance and conformance burden. Reject direct embedding for the
initial profile until it reproduces the R path; this is a scoped
engineering choice, not a claim that native integration is impossible.

The next comparison should freeze one environment and the complete
workload: data conversion, nuisance estimation, forest training,
average-effect scoring, serialization, output validation and failure
handling. Compare the maintained R route with a direct-native candidate
only under identical inputs, target, options and estimator formulas.
Exact fixtures first check coordinates, residualization, score
aggregation and refusal behavior. A separately registered generated-data
study can then examine coverage, availability, heterogeneity recovery
and resources, retaining weak-overlap and dependent cases in its
denominator. Pure arithmetic agreement cannot replace that study, and
local runtime cannot establish bank workflow savings. The useful
investment question is whether a native route improves the complete
workload enough to justify owning the additional statistical and
deployment machinery. \label{exec13:grf-adapter-end}

\section{A common interface preserves differences that matter}

The adapter comparison supports a common request and result vocabulary with deliberately different admitted operations. The Python route returns a symbolic functional under a supplied graph and explicit conditional-path restrictions. The Java route proposes a graph class under a specified discovery model and must retain failed tests before publication. The R/native route estimates a defined customer effect while preserving nuisance adjustment, aggregation and uncertainty conventions. A single effect field cannot stand in for these three services.

The local trials show why preservation extends beyond the estimator. A coordinate permutation, rounding at an integer boundary, a changed context, corrupt replay or incomplete publication can alter or conceal the scientific result. The observed refusals and successful recoveries establish particular behavior; the descriptor, input-access and error-classification findings identify repairs still needed for those versions. The next service design must carry both kinds of evidence forward. Its value is that a caller can tell what was computed, under which assumptions, what remains uncertain and which next action could resolve the limitation.

